\documentclass[11pt]{article}
\usepackage{amsmath,amssymb,amsthm}

\newtheorem{proposition}{Proposition}
\newtheorem{corollary}{Corollary}
\newtheorem{remark}{Remark}
\usepackage[sort]{natbib}
\usepackage{array}
\usepackage{booktabs}
\usepackage[flushleft]{threeparttable}
\usepackage{textcomp}
\usepackage{graphicx}
\usepackage{url}
\usepackage{multirow}
%\renewcommand{\familydefault}{\rmdefault}
\usepackage[T1]{fontenc}
%\usepackage[latin9]{inputenc}
\usepackage[a4paper]{geometry}
\geometry{verbose,tmargin=20mm,bmargin=20mm,lmargin=35mm,rmargin=20mm,footskip=0.5cm}
\pagestyle{plain}
\usepackage[usenames,dvipsnames]{xcolor}
\definecolor{dukeblue}{rgb}{0.0, 0.0, 0.61}
\definecolor{phthaloblue}{rgb}{0.0, 0.06, 0.54}
\usepackage{setspace}
\usepackage[british]{babel}
%\usepackage[sfmath]{kpfonts} %% sfmath option only to make math in sans serif. Probably only for use when the base font is sans serif.
%\renewcommand*\familydefault{\sfdefault} %% Only if the base font of the document is to be sans serif
%\usepackage[bitstream-charter]{mathdesign}
\normalfont
%     \usepackage{mathpazo}
%
\usepackage{libertine}



\title{Is the Empty Planet Unavoidable? \\
Habitat, Fertility and Long-run Growth}%\thanks{The author thanks the Monash Macroeconomics discussion group for valuable comments on the paper.}
\author{Ratbek Dzhumashev\thanks{Department of Economics, Monash University, Wellington Rd., Clayton VIC 3800, Australia and Centre for Applied Macroeconomic Analysis, Australian National University, Australia.           e-mail: \texttt{Ratbek.Dzhumashev@monash.edu}} 	
}

\date{}
\begin{document}
\maketitle
\begin{abstract}
\noindent
\citet{jones22} shows that a shrinking population can stall idea-based growth; \citet{STRULIK2024104943}
shows that growth can continue when ideas are produced by human capital.  Both study the consequences of
a declining population; neither asks what could stop the decline.  We let the taste for children and the time needed to raise them depend
on habitat per capita, which shrinks through crowding and through pollution, and let clean knowledge,
produced by publicly financed research, lower the pressure that output places on the habitat.
Replacement fertility is then a unique and locally stable outcome, and a stationary population with
positive growth is an equilibrium if and only if clean research exceeds a threshold that we derive in
closed form. When clean innovation rises with the number of people, fertility and clean technology
interact: the economies with the lowest fertility settle at the smallest populations and need the most
clean research, and below the threshold population decline slows clean innovation, so that the economy
falls into an environmental Empty Planet. Remediation of environmental quality then substitutes partly
for clean research. Fertility is positively associated with proxies for habitat per capita in four
surveys and a panel of 161 countries, and stated ideals rise with habitat at about half the rate of
realised fertility. Calibrated to these estimates, and over a range of scale effects informed by
cross-country evidence, the model implies that for a stationary equilibrium Korea needs between 1.1 and
5.4 times the clean effort that produces its current measured clean progress and Japan between 1.1 and
2.5, China and Italy fall short only at the upper end of the range, and the other five economies we
consider clear the threshold throughout. If clean knowledge is instead semi-endogenous, no research share
secures a stationary population at standard returns to research.
\end{abstract}
\onehalfspacing
\section{Introduction}
\citet{jones22} examines the economic consequences of a declining global population. In standard
idea-based growth models, sustained progress depends on a constant, or rising, pool of researchers.
As fertility in many high-income economies falls below replacement, Jones highlights the risk of an
``Empty Planet'': as population dwindles, idea production slows and living standards ultimately
stagnate. A planner could in principle avert this outcome through timely fertility policy, but
implementation delays can push the economy into a self-reinforcing trap. On this view, the long-run
fate of economic progress turns on whether societies can halt demographic decline.

Jones's pessimism follows from treating raw labour as the key input into idea creation: fewer people
mechanically imply fewer ideas. \citet{STRULIK2024104943} instead emphasises human capital as the
relevant input, so that aggregate human capital, and with it idea production, can keep expanding as
the number of workers falls. In Strulik's model with the iso-elastic ideas technology of
\citet{jones22}, the steady state with positive growth and a declining population exists but, as
\citet[Section~2.5]{STRULIK2024104943} notes, it ``is not asymptotically stable''. When idea growth
follows a bounded, logistic process, \citet[Proposition~3]{STRULIK2024104943} establishes an
asymptotically stable maximum growth rate for any fertility rate above a threshold.

Neither paper overlooks growth with a constant population. \citet[pp.~3493--3494]{jones22} recalls that
in fully endogenous growth models ``a constant population delivers constant exponential growth in living
standards forever'', and that in semi-endogenous models living standards still ``rise without bound,
albeit at a rate that slows over time''. The fertility threshold in Strulik's Proposition~3 lies below
replacement whenever $(1+\sigma)\theta>1$, where $\sigma$ is the elasticity of output with respect to
ideas and $\theta$ the elasticity of human capital with respect to education spending, so his result
applies equally to a stationary population. What
both leave open is why fertility should settle at replacement rather than below it. Strulik himself
acknowledges that ``it can be argued that steady negative population growth is also unsustainable'', and
reads his results as a guide to the next few centuries rather than to the steady state.

This paper asks whether a stationary population with sustained growth can arise as an equilibrium
outcome, and what it requires of the environment and of technology. Our answer rests on the habitat
that people share. The taste for children and the time needed to raise them depend on habitat per
capita, which falls through two channels: crowding, as more people share the same habitat, and
pollution, as production degrades the habitat itself. A declining population restores habitat per
capita and, with it, fertility, so that replacement becomes a unique and locally stable outcome rather
than a knife-edge. But output keeps growing, and the habitat shrinks unless the pressure that each unit
of output places on it falls fast enough. We make that rate endogenous. Clean knowledge, produced by
publicly financed research, lowers pressure per unit of output, and, when clean progress is proportional
to research effort, a stationary population with positive long-run growth is an equilibrium if and only if
the share of effort devoted to clean research exceeds a threshold that we derive in closed form.

The threshold links demography and technology whenever clean innovation rises with the number of people
engaged in it, as in scale-effect models of innovation. A low-fertility economy needs a large gain in
habitat per
capita before fertility returns to replacement, settles at a small population, and therefore needs more
clean research to keep pressure falling. Below the threshold the interaction is destabilising:
population decline slows clean innovation, pressure rises, habitat and fertility fall further, and the
economy converges to the rate of decline it would have without any decoupling, an environmental
counterpart of the Empty Planet of \citet{jones22}. In simulations, population undershoots on the way
to its stationary level, and an economy can fall into this trap even when the stationary state would be
sustainable once reached.

Two features of the analysis should be stated plainly. First, the long-run growth rate is set by the
ceiling of Strulik's bounded law of idea creation, which we calibrate, rather than by the demographic
mechanism. What the habitat feedback and clean technology determine is whether population stabilises,
at what level, and at what cost in clean research. Second, the bounded law is sometimes objected to
because its positive lower bound appears to permit ``innovation without people''. That bound does not
determine the long-run equilibrium here: when population stabilises at a positive level and
$(1+\sigma)\theta>1$, idea growth converges to its \emph{upper} bound. The lower bound can govern a long
transition, but not the limit.

The paper makes five contributions.
\begin{enumerate}
\item \emph{A stabilising force and endogenous decoupling.} The habitat feedback makes replacement
fertility a unique and locally stable outcome, whether it acts through the taste for children or their
cost (Proposition~\ref{prop:two_channels}). With clean progress produced by research, the condition for
a stationary population with sustained growth becomes a threshold for the clean-research share, in
closed form. Clean knowledge is assumed to earn no private return, so a stationary population is a
policy outcome; without enough clean research, the environment produces an Empty Planet of its own
(Proposition~\ref{prop:threshold}, Corollary~\ref{cor:threshold}). If clean knowledge is instead
semi-endogenous, the threshold becomes a condition on technology that no permanent research share can
meet at the returns estimated for ideas in general (Section~\ref{sub:semiendog}).
\item \emph{Fertility and clean technology interact.} With a scale effect in clean innovation, the
lowest-fertility economies need the most clean research, decline below the threshold is
self-reinforcing, and an economy can lose decoupling while its population undershoots
(Corollary~\ref{cor:threshold}, Proposition~\ref{prop:threshold}(iii), Section~\ref{sub:illustration}).
The threshold and the self-reinforcing path carry over when environmental quality acts on the cost of
children alone (Proposition~\ref{prop:split}).
\item \emph{Two instruments.} Remediation, which restores environmental quality, raises the stationary
population, one for one when quality moves tastes and costs together and by about two-fifths as much when
it moves costs alone; clean research decides whether there is a stationary population. With a scale
effect remediation also supports decoupling, but it is an expensive way to cross the threshold, and
neither instrument changes the long-run growth rate. Their value depends on how the planner values
population. An average-utilitarian planner chooses neither unless the environment is valued directly
and, for remediation, the tax spares education spending. A
total-utilitarian planner, at our baseline, chooses both, and enough clean research to cross the
threshold, except where fertility is lowest and the scale effect strong: in Korea, no criterion we
consider favours crossing it (Section~\ref{sec:policy}, Appendix~\ref{sub:welfare}).
\item \emph{Evidence.} Fertility is positively associated with proxies for habitat per capita in four
surveys and a panel of 161 countries, and, under the maintained conversion from settlement size to
habitat, the implied feedback strength, $\hat S\approx0.21$, lies well inside the stable region. Stated
ideal fertility rises with habitat at a little over half the rate of realised fertility, pointing to a
role for preferences alongside costs. At a given density and income, fertility is higher where the
pollution load is lower, while stated ideals show no detectable response to pollution, suggesting that
pollution acts mainly through costs. The evidence is associational: it documents the correlations the mechanism
predicts, not causal effects (Section~\ref{sec:empirics}).
\item \emph{Quantitative implications.} Calibrated to these estimates, and with clean progress read from
each country's observed decline in emission intensity, the model implies that for a stationary
equilibrium Korea needs between 1.1 and 5.4 times the clean effort that produces its current measured
clean progress and Japan between 1.1 and 2.5, over a sensitivity range for the scale effect,
$0\le\nu\le0.5$, that is informed by cross-country evidence. China and Italy clear the threshold
without a scale effect and fall short at the upper end of the range, by factors of 1.8 and 1.4, while
Germany, the United States, France, Sweden and the United Kingdom clear it throughout. Starting from
today's conditions, transitions add up to a tenth to these requirements. Population stabilises only
after large declines, and the lower the fertility, the larger the decline: at the within-country
pollution elasticity, to about half its present level in France and the United States and to a few per
cent in Korea, although a larger pollution elasticity would raise these levels substantially
(Section~\ref{sec:quant}).
\end{enumerate}

\paragraph{Related literature.}
Two recent contributions are closest to ours. \citet{beulmann2025} study how fertility affects
emissions and green innovation, and \citet{edenkuruc2026} how the scale of a stationary population
affects income. Neither asks whether an environmental feedback stops population decline, which is the
question of this paper.

More broadly, the paper contributes to three strands of work. The first links population, ideas and growth, from
\citet{romer1990} and \citet{kremer1993} to the semi-endogenous models of \citet{jones1995}, with
\citet{bloom2020} documenting that ideas are getting harder to find. \citet{peretto2015} study growth on
a finite planet in which resource scarcity and population interact. \citet{edenkuruc2026} take
stabilisation as given, ``perhaps via \ldots\ policy that induces replacement fertility'', and show that a
larger stationary population is likely to be richer per head, because the innovation benefits of scale
outweigh the dilution of natural resources. Relative to this strand, we supply what is usually assumed: a
reason why population stabilises, and a determination of the level at which it does.

The second strand concerns the determinants of fertility. Our household problem descends from
\citet{barro1989} and the quantity--quality trade-off of \citet{beckerlewis1973} and of unified growth
theory \citep{galor2000,doep05}. Within this tradition the cost of children is itself endogenous: parents
can substitute purchased childcare for their own time, so that the effective cost of a child depends on
wages and on the price of childcare \citep{hazan15,day2016}, and fertility responds to social norms and to
the consumption and fertility of others \citep{palivos2001,munshi2006}. \citet{Dz} combine endogenous
childcare costs with such social externalities and show that together they generate both the fertility
decline of the demographic transition and the subsequent rebound in high-income countries
\citep{myrskyla2009}; \citet{jones11} review how far these theories account for the negative
fertility--income relationship. The idea that crowding lowers fertility by raising the cost of children
goes back to \citet{goodsell} and \citet{thompson}. \citet{DAO2013442} use it to show how population density
can lock a society in stagnation, and \citet{DELACROIX201713} document the negative association between
density and fertility in developing countries. \citet{dettling2014} show that house prices affect
birth rates, a goods-cost channel that Corollary~\ref{cor:goods} adds to the model and
Appendix~\ref{app:budget} takes to the data. Outside economics, \citet{lutz2006} argue that density is a
key factor in declining fertility, and \citet{sng2017} that crowding shifts people towards slower life
strategies with fewer, more heavily invested children. On the environmental side,
\citet{dao25} show that climate-induced depletion of local resources raises the time women spend collecting
water and firewood, widens the gender gap in education and slows the fertility decline, while crowding
raises the time cost of children and stabilises population. In their model environmental degradation
raises fertility, whereas in ours a degraded habitat lowers it; crowding lowers fertility in both.
\citet{mari10}, \citet{Varvarigos2013} and
\citet{schaefer2020} let pollution affect fertility through health and child survival. We do not model
mortality explicitly, and our habitat channel can be read as a reduced form that absorbs such effects.
Relative to this
strand, we make the taste for children, and not only their cost, depend on a habitat that is shared and
endogenous, and embed the mechanism in an idea-based growth model.

The third strand examines how production affects the environment and how the environment feeds back on
population. In \citet{raffin14}, \citet{econ2008} and \citet{marsiglio2018}, output growth can improve
environmental quality when abatement is strong enough, whereas in \citet{Varvarigos2013} and
\citet{pautrel08} production growth degrades it despite abatement; \citet{mar2017} lets both output and
population drive emissions. \citet{DZHUMASHEV2024} shows that physical output, because it has mass, cannot
be produced without material inputs: the marginal product of resources is bounded, so that the Inada
condition for resources fails and growing output requires growing resource use and, potentially, more
pollution \citep[see also][]{baum04}. \citet{dzhumashev2025} show that this resource use degrades the
environment and lowers fertility, creating a trade-off between income growth and population growth.
In the resource-based models of \citet{Brander1998} and \citet{lagerlof2006}, population responds to
the resource stock strongly enough to oscillate.
\citet{DELACROIX2012271} show that pollution control, by taxing production, shifts activity towards
untaxed procreation and can impoverish later generations, and \citet{gerlagh2023} combine a carbon tax
with a tax on births. In \citet{beulmann2025}, clean technology results from
directed technical change: lower fertility raises human capital, output and emissions at first, and
speeds the transition to green growth once research is redirected. We study the reverse direction, from
the environment to fertility, and make clean technology a product of publicly financed research, so that
fertility and clean innovation interact. Relative to this strand, we derive the clean-research threshold
at which the loop delivers a stationary population with growth, show how it depends on fertility, and
compare clean research with remediation.

The rest of the paper is organised as follows. Section~\ref{sec:model} sets out the model.
Section~\ref{sec:theory} derives the stationary habitat, the clean-research threshold and its
implications, illustrates them numerically, and extends them to semi-endogenous clean knowledge and to
quality that acts on the costs of children only. Section~\ref{sec:empirics} presents the evidence on the
habitat--fertility link, and Section~\ref{sec:quant} calibrates the habitat block to it and computes the
clean-research thresholds of nine countries. Section~\ref{sec:policy} discusses policy and
Section~\ref{sec:conclusion} concludes. The appendices report the budget-based tests, a cross-country
consistency check, the welfare analysis of remediation and clean research, simulated country paths, and
cross-country evidence on the scale effect.



\section{The model}
\label{sec:model}
\label{sec:habitat_taste}

\subsection{Demography and habitat per capita}
Time is discrete, $t=0,1,2,\dots$. In each period a generation of homogeneous adults makes fertility
and education decisions for their children, who become workers in period $t+1$. There is no physical
capital. Let $L_t$ denote the adult population, $n_t$ fertility, $h_t$ human capital per worker and
$\ell_t$ labour supply per worker. Effective labour is $H_t=h_t\ell_tL_t$, and population evolves as
\begin{align}
L_{t+1}=n_tL_t .
\label{eq:Llaw_hab}
\end{align}
Fertility $n_t$ is the generational reproduction factor, normalised so that $n_t=1$ is replacement.
Because each child has two parents, it equals half the number of surviving children per woman. When we
take the model to data we convert observed fertility in the same way: children ever born are divided by
two, and the total fertility rate by the replacement rate of $2.1$, which also allows for mortality
before adulthood.

Let $E_t$ denote the effective human habitat, the land available to the population weighted by its
environmental quality, $E_t=\bar E\,Q_t$, where $\bar E$ is the fixed land area and $Q_t>0$ an index of
quality. Habitat per capita is
\begin{align}
\epsilon_t\equiv\frac{E_t}{L_t}=\frac{Q_t}{D_t},
\qquad D_t\equiv\frac{L_t}{\bar E}.
\label{eq:eps_def}
\end{align}
It falls for two distinct reasons: crowding, a higher density $D_t$, and pollution, a lower quality
$Q_t$. Habitat per capita enters the household problem in two places. The taste for children rises
with it, and the time needed to raise a child falls with it:
\begin{align}
\alpha_t=\alpha(\epsilon_t),\quad \alpha'(\epsilon)\ge0,
\qquad
\phi_t=\phi(\epsilon_t),\quad \phi'(\epsilon)\le0,\quad 0<\phi(\epsilon)<1,
\label{eq:alpha_eps}
\end{align}
with at least one inequality strict. A crowded or degraded habitat may lower the desire for children,
as in the literature on density and life strategies \citep{lutz2006,sng2017}, or raise the time that children demand, through crowding \citep{DAO2013442,dao25} and through child
health \citep{mari10,Varvarigos2013}. We keep both channels
because the evidence of Section~\ref{sec:empirics} points to both. Section~\ref{sub:split} lets crowding
and environmental quality act through different channels. Households take $\epsilon_t$, and hence
$\alpha_t$ and $\phi_t$, as given.

\subsection{Production, ideas and clean knowledge}
A share $R_t$ of effective labour works in clean research, described below, and the rest in production:
\begin{align}
Y_t=A_t^{\sigma}(1-R_t)H_t,\qquad \sigma\ge0,
\label{eq:Y_hab}
\end{align}
so that the wage per unit of human capital is $w_t=A_t^{\sigma}$. The stock of ideas follows the bounded
law of \citet{STRULIK2024104943}, driven by effective labour in production,
\begin{align}
A_{t+1}=(1+g_{A,t})A_t,
\qquad
g_{A,t}=\frac{\bar g_A}{1+\omega\exp\!\left(-(1-R_t)H_t/A_t\right)},
\qquad \bar g_A>0,\ \omega>0,
\label{eq:gA_logistic}
\end{align}
which, as Strulik notes, implements a learning-by-doing assumption. Idea growth is bounded above,
so that it saturates, and below by $\underline g_A=\bar g_A/(1+\omega)>0$.\footnote{With the
iso-elastic law of \citet{jones22}, $A_{t+1}-A_t=A_t^{1-\beta}H_t^{\lambda}$, the steady state with a
declining population exists but, as \citet[Section~2.5]{STRULIK2024104943} shows, is not
asymptotically stable. We use the bounded law throughout.}

Production degrades the habitat. Because output has mass, producing more of it requires more material
input, and the marginal product of resources is bounded \citep{DZHUMASHEV2024}; output cannot grow
without a growing flow of materials, and with it of waste and emissions, unless technology reduces the
pressure that each unit of output places on the environment. We capture that possibility with a stock of
clean knowledge $B_t$ and define pressure on the habitat as output per unit of clean knowledge,
\begin{align}
P_t=\frac{Y_t}{B_t},
\qquad
B_{t+1}=\left(1+\delta_B R_tL_t^{\nu}\right)B_t,
\qquad \delta_B>0,\ 0\le\nu\le1 .
\label{eq:B_law}
\end{align}
Clean progress depends on the effort devoted to it, $R_t$, and, when $\nu>0$, on the size of the
population, and hence on the number of clean researchers. With $\nu=0$ it depends on research intensity alone, as in
growth models without scale effects; with $\nu=1$ it is proportional to the number of clean
researchers, $R_tL_t$, as in the scale-effect models of \citet{romer1990} and \citet{jones1995}. Because
adults within a generation share the same human capital, $R_t$ is both the share of effective labour
and the share of workers in clean research, and the scale effect counts researchers by headcount.
A common shortcut lets pressure per unit of output fall at an exogenous rate tied to knowledge,
$P_t=Y_t/A_t^{\zeta}$. With $\nu=0$ and constant $R$, \eqref{eq:B_law} is equivalent to it in the
saturated limit, with $\zeta=\ln(1+\delta_BR)/\ln(1+\bar g_A)$. Clean knowledge lowers an externality and earns no private return, so that without policy
$R_t=0$; this is the micro-foundation we leave aside, and \citet{beulmann2025} model it with directed
technical change.

Habitat depends on pressure and on remediation,
\begin{align}
E_t=\bar E\,\mathcal{Q}(P_t;\tau_r),
\qquad \mathcal{Q}_P<0<\mathcal{Q}_\tau,
\qquad \mathcal{Q}(0;\tau_r)=Q_\infty(\tau_r),
\label{eq:pressure}
\end{align}
where $\mathcal{Q}(P;\tau_r)$ behaves as $P^{-\mu}$, $\mu>0$, for large $P$, and $Q_\infty(\tau_r)$ is the
quality of a habitat under no pressure. The government taxes labour income at the rate $\tau_r+R$. It
spends the revenue $\tau_r w_tH_t$ on remediation, which restores quality at given pressure, and with the
revenue $Rw_tH_t$ it hires the share $R$ of effective labour at the market wage for clean research. The
clean-research share therefore equals its tax rate, and both instruments are financed in the same way.

\subsection{Households}
Adults value consumption $c_t$, the number of their children $b_t$ and the human capital of their
children $h_{t+1}$:
\begin{align}
u_t=\log c_t+\alpha(\epsilon_t)\log b_t+\gamma\log h_{t+1},\qquad b_t\equiv2n_t,\qquad \gamma>0 .
\label{eq:util_hab}
\end{align}
Children are counted per couple, $b_t$, while the population law \eqref{eq:Llaw_hab} uses $n_t$. Since
$\log b_t=\log2+\log n_t$, the factor two leaves every choice unchanged; it matters only for welfare.
Raising a child takes $\phi(\epsilon_t)$ units of adult time; parents also spend $s_t$ units of time and
$q_t$ units of goods per child on education, and human capital is produced as
\begin{align}
h_{t+1}=\nu_h\,q_t^{\theta}s_t^{\eta},\qquad \nu_h>0,\ 0<\theta<1,\ 0<\eta<1 .
\label{eq:hc_hab}
\end{align}
With $\tau\equiv\tau_r+R$ the budget constraint is
\begin{align}
(1-\tau)\,w_th_t\left(1-\phi(\epsilon_t)n_t-s_tn_t\right)=c_t+q_tn_t .
\label{eq:budget_hab}
\end{align}

\begin{proposition}[Household choices]
\label{prop:hab_solution}
Let $\Gamma\equiv\gamma(\eta+\theta)$, write $\alpha_t=\alpha(\epsilon_t)$ and $\phi_t=\phi(\epsilon_t)$, and
assume $\alpha_t>\Gamma$. Then
\begin{align}
c_t&=\frac{(1-\tau)w_th_t}{1+\alpha_t},
\qquad
n_t=n(\alpha_t,\phi_t)\equiv\frac{\alpha_t-\Gamma}{\phi_t(1+\alpha_t)},
\label{eq:sol_n_hab}\\
s_t&=\frac{\eta\gamma\phi_t}{\alpha_t-\Gamma},
\qquad
q_t=q^\ast(\alpha_t,\phi_t)\,(1-\tau)w_th_t,\quad q^\ast\equiv\frac{\gamma\theta\phi_t}{\alpha_t-\Gamma},
\label{eq:sol_q_hab}\\
\ell_t&=\ell(\alpha_t)\equiv\frac{1+\gamma\theta}{1+\alpha_t}.
\label{eq:sol_l_hab}
\end{align}
\end{proposition}

\begin{proof}
With multiplier $\Omega$ on \eqref{eq:budget_hab} and $W_t\equiv(1-\tau)w_th_t$, the first-order
conditions for $c_t$, $q_t$ and $s_t$ give $q_tn_t=\gamma\theta c_t$ and $W_ts_tn_t=\gamma\eta c_t$. The
condition for $n_t$, $\alpha_t/n_t=\Omega[W_t(\phi_t+s_t)+q_t]$, becomes
$(\alpha_t-\Gamma)c_t=\phi_tW_tn_t$. Substituting the three relations into the budget constraint gives
$c_t(1+\Gamma)=W_t-(\alpha_t-\Gamma)c_t$, hence $c_t=W_t/(1+\alpha_t)$, and the rest follows, labour
supply from $\ell_t=1-\phi_tn_t-s_tn_t$.
\end{proof}

Two features carry most of what follows. Fertility responds to habitat through both channels, and its
habitat elasticity is
\begin{align}
\frac{\mathrm{d}\ln n}{\mathrm{d}\ln\epsilon}=\mathcal{K}(\alpha)\,\xi+\chi,
\qquad
\mathcal{K}(\alpha)\equiv\frac{\alpha(1+\Gamma)}{(\alpha-\Gamma)(1+\alpha)},
\qquad
\xi\equiv\frac{\mathrm{d}\ln\alpha}{\mathrm{d}\ln\epsilon},\quad
\chi\equiv-\frac{\mathrm{d}\ln\phi}{\mathrm{d}\ln\epsilon}.
\label{eq:S_def}
\end{align}
Labour supply, by contrast, contains neither $\phi$ nor the taxes: a cheaper child is met by more
children and less education per child, leaving time at work unchanged. Taxes scale the return to all
labour time and so leave every time allocation unchanged.

\begin{corollary}[Goods costs]
\label{cor:goods}
If each child also requires goods worth $\kappa_t=k(\epsilon_t)W_t$ with $k'(\epsilon)\le0$, such as
housing, fertility is $n_t=(\alpha_t-\Gamma)/[(1+\alpha_t)\Phi(\epsilon_t)]$ with
$\Phi\equiv\phi+k$, the full cost of a child as a share of full income. The results on fertility, the
stationary habitat and the clean-research threshold hold with $\phi$ replaced by $\Phi$ and $\chi$ by
$\chi_\Phi\equiv-\mathrm{d}\ln\Phi/\mathrm{d}\ln\epsilon$; labour supply then also depends on $\phi/\Phi$
(Appendix~\ref{app:budget}).
\end{corollary}

A goods cost, unlike a time cost, moves labour supply, and with it the fixed point
\eqref{eq:alpha_fixed_point} below; Appendix~\ref{app:budget} uses the distinction,
and the invariance of expenditure shares to every cost, to separate tastes from costs in budget data.

\subsection{Within-period equilibrium}
Output depends on the taste for children through labour supply, and habitat depends on output, so the
taste for children solves the fixed point
\begin{align}
\alpha_t=\alpha\!\left(\frac{\bar E\,\mathcal{Q}\big(A_t^{\sigma}(1-R)h_t\ell(\alpha_t)L_t/B_t;\tau_r\big)}{L_t}\right).
\label{eq:alpha_fixed_point}
\end{align}
Because $\ell$ does not depend on $\phi$, the time cost does not enter \eqref{eq:alpha_fixed_point}; it is
determined residually as $\phi_t=\phi(\epsilon_t)$. Since $\ell'<0$ and $\mathcal{Q}_P<0$, the right-hand
side is increasing in $\alpha_t$, and we maintain the condition that its slope is below one, which makes
the fixed point locally unique and ensures that
$J_t\equiv1-\alpha'(\epsilon_t)\,\bar E\,\mathcal{Q}_P\,A_t^{\sigma}(1-R)h_t\ell'(\alpha_t)/B_t>0$.

\begin{proposition}[Within-period effects]
\label{prop:tax}
\begin{enumerate}
\item[(i)] A more productive, better educated or larger population crowds the habitat and lowers
habitat per capita and fertility, and does not raise the taste for children; a larger stock of clean
knowledge has the opposite effects.
\item[(ii)] Remediation raises habitat per capita and fertility:
\begin{align}
\frac{\partial\ln n_t}{\partial\tau_r}=\big[\mathcal{K}(\alpha_t)\,\xi+\chi\big]\,
\frac{\partial\ln\epsilon_t}{\partial\tau_r},
\qquad
\frac{\partial\ln\epsilon_t}{\partial\tau_r}=\frac{\mathcal{Q}_\tau}{\mathcal{Q}\,J_t}>0 .
\label{eq:dn_dtau}
\end{align}
\item[(iii)] Holding $A_t$ and $h_t$ fixed, clean research lowers income per capita, and remediation
does not raise it, lowering it when $\alpha'>0$ through labour supply. Both taxes lower the human capital of
the next generation, which on the stationary path scales with $(1-\tau_r-R)^{\theta/(1-\theta)}$.
\end{enumerate}
\end{proposition}

\begin{proof}
Parts (i) and (ii) follow from the implicit function theorem applied to \eqref{eq:alpha_fixed_point}
and from \eqref{eq:S_def}. For (iii), $y_t=A_t^{\sigma}(1-R)h_t\ell(\alpha_t)$, where remediation
raises $\alpha_t$ when $\alpha'>0$ and leaves it unchanged otherwise, and $\ell'<0$; and
$h_{t+1}=\nu_h\big((1-\tau)q^\ast A_t^{\sigma}h_t\big)^{\theta}s_t^{\eta}$; with $\alpha$ and $\phi$ constant,
$\ln h$ follows a stable first-order difference equation whose intercept contains
$\theta\ln(1-\tau)$, which gives the long-run factor.
\end{proof}

Given $\alpha_t$, the economy evolves according to \eqref{eq:Llaw_hab}, \eqref{eq:gA_logistic},
\eqref{eq:B_law} and
\begin{align}
h_{t+1}=\nu_h\big((1-\tau)\,q^\ast(\alpha_t,\phi_t)\,A_t^{\sigma}h_t\big)^{\theta}s_t^{\eta}.
\label{eq:h_tax_dyn}
\end{align}

\section{Stationary population and endogenous decoupling}
\label{sec:theory}
\label{sec:timecost}
\label{sec:endog_env}

\subsection{The stationary habitat}
\label{sub:stationary}
Suppose first that the habitat is fixed. Then $\epsilon_{t+1}=f(\epsilon_t)\equiv\epsilon_t/n(\epsilon_t)$,
and a stationary point requires replacement fertility,
\begin{align}
\alpha(\epsilon^\ast)=\frac{\Gamma+\phi(\epsilon^\ast)}{1-\phi(\epsilon^\ast)} .
\label{eq:stationary_phi}
\end{align}

\begin{proposition}[The stationary habitat]
\label{prop:two_channels}
\label{rem:eps_stability_hab}
Suppose $\alpha(\epsilon)-[\Gamma+\phi(\epsilon)]/[1-\phi(\epsilon)]$ changes sign on $(0,\infty)$. Then
\begin{enumerate}
\item[(i)] the stationary habitat per capita $\epsilon^\ast$ exists and is unique;
\item[(ii)] the strength of the demographic feedback at $\epsilon^\ast$ is
$S\equiv\epsilon^\ast n'(\epsilon^\ast)=\mathcal{K}(\alpha^\ast)\xi+\chi$, with $\alpha^\ast=\alpha(\epsilon^\ast)$;
\item[(iii)] $\epsilon^\ast$ is locally stable if and only if $0<S<2$; convergence is monotone for
$0<S<1$ and takes the form of damped oscillations in fertility and population for $1<S<2$.
\end{enumerate}
\end{proposition}

\begin{proof}
The left-hand side of \eqref{eq:stationary_phi} is non-decreasing in $\epsilon$ and the right-hand side
non-increasing, with one of the two strict, so their difference crosses zero once. Part (ii) is
\eqref{eq:S_def} at $n=1$, and (iii) follows from $f'(\epsilon^\ast)=1-S$.
\end{proof}

Population stabilisation is therefore not hostage to the taste channel: either a taste or a cost that
responds to habitat suffices. The feedback strength $S$ governs the speed of adjustment, with a
half-life of $\ln0.5/\ln|1-S|$ generations, and it is bounded above as well as below: too strong a
response makes fertility oscillate, as in the resource-based models of \citet{Brander1998} and
\citet{lagerlof2006}, and eventually destabilises the stationary point. Our estimate in
Section~\ref{sub:xi}, $\hat S\approx0.21$, places the economy in the monotone region, with a half-life of
about three generations.

\subsection{Saturation of idea growth}
With a positive stationary population, what matters for the ideas law is the ratio of effective labour
to knowledge. Along a path with constant growth rates, \eqref{eq:h_tax_dyn} implies
$1+g_h=(1+g_A)^{\sigma\theta/(1-\theta)}$, so that $H_t/A_t$ drifts at the rate
$(1+g_A)^{\sigma\theta/(1-\theta)-1}$.

\begin{proposition}[Saturation]
\label{prop:saturation}
Let $L_t\to L^\ast>0$ and $\alpha_t\to\alpha^\ast>\Gamma$. If $(1+\sigma)\theta>1$, then $H_t/A_t\to\infty$
and $g_{A,t}\to\bar g_A$, and income per capita grows at
\begin{align}
1+g_y=(1+\bar g_A)^{\zeta^\ast},
\qquad \zeta^\ast\equiv\frac{\sigma}{1-\theta}.
\label{eq:limits_saturated}
\end{align}
If $(1+\sigma)\theta<1$, idea growth converges to its lower bound $\underline g_A$.
\end{proposition}

\begin{proof}
Write $a_t=\ln A_t$, $b_t=\ln h_t$ and $z_t=b_t-a_t$. Then
$z_{t+1}=\text{const}+\theta z_t+[(1+\sigma)\theta-1]a_t-\ln(1+g_{A,t})$. Since $g_{A,t}\ge\underline g_A>0$,
$a_t\to\infty$, and since $0<\theta<1$, the sign of $(1+\sigma)\theta-1$ decides whether $z_t$ diverges to
plus or minus infinity. Applying \eqref{eq:gA_logistic} gives the limits.
\end{proof}

\begin{remark}[The lower bound governs transitions]
\label{rem:transition}
Proposition~\ref{prop:saturation} is asymptotic. Along the transition $H_t/A_t$ may first fall for many
generations, so that $g_{A,t}$ stays close to its lower bound for a long time; the upper bound governs
the limit. Long-run growth is thus set by the ceiling $\bar g_A$, which we calibrate, and is the same
whether population is stationary or declines slowly enough for idea growth to saturate
(Proposition~\ref{prop:threshold}). What the habitat mechanism determines is whether
population settles, at what level, and what that requires of clean technology.
\end{remark}

\subsection{Endogenous decoupling}
\label{sub:decoupling}
With output growing without bound, pressure stays bounded only if clean knowledge grows at least as fast
as output, and vanishes only if it grows faster. At a stationary population, output grows by $(1+\bar g_A)^{\zeta^\ast}$ per generation, which
defines the threshold for clean progress,
\begin{align}
g_B^\ast\equiv(1+\bar g_A)^{\zeta^\ast}-1,
\label{eq:gB_star}
\end{align}
and the stationary population that a habitat under no pressure supports,
\begin{align}
L^\ast(\tau_r)\equiv\frac{\bar E\,Q_\infty(\tau_r)}{\epsilon^\ast}.
\label{eq:L_star}
\end{align}

\begin{proposition}[Endogenous decoupling and the stationary population]
\label{prop:threshold}
\label{prop:decoupling}
Let $(1+\sigma)\theta>1$, let $\epsilon^\ast$ satisfy $0<S<2$, and let $R\in(0,1)$ and $\tau_r$ be constant.
\begin{enumerate}
\item[(i)] If $\delta_BR\,L^\ast(\tau_r)^{\nu}>g_B^\ast$, the stationary population $L^\ast(\tau_r)$, with
income per capita growing at the saturated rate \eqref{eq:limits_saturated}, is a locally stable limit:
pressure vanishes, habitat tends to $\bar EQ_\infty(\tau_r)$, habitat per capita to $\epsilon^\ast$ and
fertility to replacement.
\item[(ii)] If $\delta_BR\,L^\ast(\tau_r)^{\nu}<g_B^\ast$, no stationary population is compatible with
growing output. With $\nu=0$, population declines asymptotically at
\begin{align}
\tilde n(R)=(1+\bar g_A)^{\frac{\mu\,(\zeta(R)-\zeta^\ast)}{1+\mu}},
\qquad
\zeta(R)\equiv\frac{\ln(1+\delta_BR)}{\ln(1+\bar g_A)}.
\label{eq:n_tilde}
\end{align}
With $\nu>0$, clean progress vanishes along any declining path and population declines at the rate
without any decoupling, $\tilde n_0=(1+\bar g_A)^{-\mu\zeta^\ast/(1+\mu)}$.
\item[(iii)] With $\nu>0$, the declining path of (ii) is self-consistent also when the condition in (i)
holds: along it population falls, clean progress $\delta_BRL_t^{\nu}$ shrinks and pressure keeps rising,
so the path never reaches the stationary population.
\end{enumerate}
Throughout, fertility stays in the interior range of Proposition~\ref{prop:hab_solution},
$\Gamma<\alpha_t$ and $n_t<1/\phi_t$, and the declining paths require $\tilde n$ above
$(1+\bar g_A)^{(1-(1+\sigma)\theta)/(1-\theta)}$, the threshold of
\citet[Proposition~3]{STRULIK2024104943} above which idea growth saturates along a declining population.
Part (iii) establishes that the declining path exists alongside the stationary one, not which initial
conditions lead to it; Section~\ref{sub:illustration} shows by simulation that an undershoot of
population on the way to $L^\ast$ can.
\end{proposition}

\begin{proof}
At the saturated limit $h_t\propto A_t^{\sigma\theta/(1-\theta)}$, so output grows by
$(1+\bar g_A)^{\zeta^\ast}n_t$ per generation and pressure by
$(1+\bar g_A)^{\zeta^\ast}n_t/(1+\delta_BRL_t^{\nu})$. (i) Near $L^\ast$ the strict inequality keeps
pressure falling, so $P_t\to0$, $E_t\to\bar EQ_\infty(\tau_r)$, $E_{t+1}/E_t\to1$, and habitat per capita
follows the fixed-habitat law, whose stationary point $\epsilon^\ast$ is locally stable by
Proposition~\ref{prop:two_channels}. (ii) A stationary population requires bounded pressure and hence
$\delta_BRL^{\nu}\ge g_B^\ast$ at a population no larger than $L^\ast(\tau_r)$, which the inequality
excludes. On a path with constant habitat per capita, $E$ and $L$ shrink at a common rate $\tilde n$,
and for large pressure $E_{t+1}/E_t=(P_{t+1}/P_t)^{-\mu}$. With $\nu=0$, setting
$[(1+\bar g_A)^{\zeta^\ast}\tilde n/(1+\delta_BR)]^{-\mu}=\tilde n$ gives \eqref{eq:n_tilde}. With
$\nu>0$, $L_t\to0$ implies $\delta_BRL_t^{\nu}\to0$, and the same calculation with $\zeta=0$ gives
$\tilde n_0$. (iii) Along the declining path of (ii), $L_t\to0$ and hence $\delta_BRL_t^{\nu}\to0$ whatever
the value of $L^\ast(\tau_r)$, so the path satisfies the conditions used in (ii) at every date and is an
equilibrium path. The argument establishes self-consistency only; we do not characterise its basin of
attraction.
\end{proof}

The proposition replaces an assumption by a result. Under exogenous decoupling, a fixed rate $\zeta$
decides whether population settles, and policy cannot affect it by construction. Here the
rate of clean progress is chosen through $R$, and whether it suffices depends on demography.

\begin{corollary}[The clean-research threshold]
\label{cor:threshold}
The smallest clean-research share at which the stationary population is a locally stable equilibrium is
\begin{align}
R^\ast(\tau_r)=\frac{g_B^\ast}{\delta_B}\left(\frac{\epsilon^\ast}{\bar E\,Q_\infty(\tau_r)}\right)^{\nu}.
\label{eq:R_star}
\end{align}
\begin{enumerate}
\item[(a)] \emph{Remediation substitutes for clean research if and only if $\nu>0$:}
\begin{align*}
\frac{\partial\ln R^\ast}{\partial\tau_r}=-\nu\,\frac{Q_\infty'(\tau_r)}{Q_\infty(\tau_r)}.
\end{align*}
\item[(b)] \emph{Low fertility raises the threshold:} $\partial\ln R^\ast/\partial\ln\epsilon^\ast=\nu$. An
economy whose taste for children is low, or whose cost of children is high, needs a larger habitat per
capita to restore replacement, settles at a smaller population, and needs more clean research.
\item[(c)] \emph{Neither instrument changes long-run growth.} Both lower the level of income per capita,
through human capital as in Proposition~\ref{prop:tax}(iii), and clean research also by the factor $1-R$
in production; both leave the saturated growth rate \eqref{eq:limits_saturated} unchanged.
\end{enumerate}
\end{corollary}

\begin{proof}
Equation \eqref{eq:R_star} solves $\delta_BR(L^\ast)^{\nu}=g_B^\ast$ with \eqref{eq:L_star}; (a) and (b)
follow by differentiation, and (c) from Propositions~\ref{prop:tax} and~\ref{prop:saturation}, since the
taxes enter \eqref{eq:h_tax_dyn} only through the constant factor $(1-\tau)^{\theta}$.
\end{proof}

The threshold $R^\ast$ is a condition for the stationary equilibrium to exist and attract nearby paths.
By Proposition~\ref{prop:threshold}(iii), an economy that starts far from $L^\ast$ can need more clean
research to get there; Section~\ref{sub:illustration} measures the difference, which we call the
transition premium. Three features of the corollary deserve emphasis. First, because clean knowledge earns no private
return, the market supplies no clean research, so a stationary population is a policy outcome. Second, without a scale effect, remediation and clean
research are not substitutes: remediation buys a larger stationary population, while only clean research
decides whether there is one. With a scale effect, remediation also supports decoupling, because the
larger population it sustains does more clean research. It is nevertheless an expensive way to cross the
threshold: the total tax $\tau_r+R^\ast(\tau_r)$ falls with $\tau_r$ only if $\nu\,R^\ast\,Q_\infty'/Q_\infty>1$,
which fails when the threshold share is a few per cent. Third, fertility and clean technology interact in
both directions. Low fertility raises the clean effort required, and a shortfall of clean effort lowers
fertility; with $\nu>0$ the second link makes decline self-reinforcing, an environmental counterpart of
the Empty Planet of \citet{jones22}.

\subsection{Numerical illustration}
\label{sub:illustration}
\label{sec:calibration}
We adopt the calibration of \citet[Section~3.2]{STRULIK2024104943} for preferences, education and
ideas: $\alpha_0=0.39$, $\gamma=0.26$, $\eta=0.07$, $\phi_0=0.138$, $\theta=0.83$, $\sigma=1$, $\omega=100$
and $\nu_h=11$, so that $\Gamma=0.234$, $\alpha^\ast=(\phi_0+\Gamma)/(1-\phi_0)=0.432$ and fertility
at $\alpha_0$ is $0.813$, a population decline of $0.69$ per cent a year with thirty-year generations.
We set $\bar g_A=0.0975$ per generation, so that income per capita grows at $1.84$ per cent a year in the
saturated limit, close to Strulik's benchmark. For this illustration the habitat acts through tastes
alone, $\alpha(\epsilon)=\alpha_0\epsilon^{\xi}$, with $\xi$ chosen so that doubling habitat per capita
restores replacement, $\xi=\ln(\alpha^\ast/\alpha_0)/\ln2=0.146$, which gives
$\mathcal{K}(\alpha^\ast)=1.883$ and $S=0.275$. Habitat is $\bar E(1+r\tau_r)/(1+P/\bar P)$, so that
$\mu=1$, with $\bar E=2$, $r=2$ and $\bar P$ set so that habitat per capita starts at one; the stationary
population is then $L^\ast=1+2\tau_r$ in units of the initial one. We scale clean productivity so that
$R=2$ per cent gives clean knowledge the growth that decoupling at $\zeta=7$ would give; only the
ratio $R/R^\ast$ matters for what follows. Section~\ref{sec:quant} replaces these illustrative values with
estimates.

\begin{table}[h!]\centering
\begin{threeparttable}
\caption{The clean-research threshold: closed form and full-model simulation}
\label{tab:threshold}
\footnotesize
\begin{tabular}{lcccccc}
\toprule
 & \multicolumn{2}{c}{$\nu=0$} & \multicolumn{2}{c}{$\nu=0.5$} & \multicolumn{2}{c}{$\nu=1$}\\
\cmidrule(lr){2-3}\cmidrule(lr){4-5}\cmidrule(lr){6-7}
 & closed form & simulated & closed form & simulated & closed form & simulated\\
\midrule
\multicolumn{7}{l}{\emph{A. Remediation, $\alpha_0=0.39$ ($\epsilon^\ast=2$)}}\\
$\tau_r=0$   & 1.59 & 1.58 & 1.59 & 1.66 & 1.59 & 1.80\\
$\tau_r=0.1$ & 1.59 & 1.57 & 1.45 & 1.47 & 1.32 & 1.42\\
$\tau_r=0.2$ & 1.59 & 1.56 & 1.34 & 1.33 & 1.13 & 1.15\\
\midrule
\multicolumn{7}{l}{\emph{B. Lower taste for children, $\tau_r=0$}}\\
$\alpha_0=0.36$ ($\epsilon^\ast=3.5$)  & 1.59 & & 2.09 & 2.10 & 2.75 & \\
$\alpha_0=0.33$ ($\epsilon^\ast=6.3$)  & 1.59 & & 2.81 & 2.76 & 4.99 & \\
$\alpha_0=0.30$ ($\epsilon^\ast=12.1$) & 1.59 & & 3.90 & 3.67 & 9.56 & \\
\bottomrule
\end{tabular}
\begin{tablenotes}\footnotesize
\item Clean-research share, per cent of effective labour. Closed form: \eqref{eq:R_star}. Simulated: the
smallest share at which the full model, with the within-period fixed point solved each generation,
settles at replacement within 400 generations. Simulated values below the closed form, by up to 6 per
cent in panel B, arise because, just below the threshold, pressure grows so slowly that the decline does
not show within the horizon.
\end{tablenotes}
\end{threeparttable}
\end{table}

\begin{figure}[h!]\centering
\includegraphics[width=\textwidth]{fig_threshold.pdf}
\caption{The clean-research threshold. Panel (a): threshold against the remediation tax, closed form
\eqref{eq:R_star} (lines) and full-model simulation (markers). Panel (b): population paths with $\nu=0.5$
and $\tau_r=0$ for clean research just above and below the simulated threshold.}
\label{fig:threshold}
\end{figure}

Table~\ref{tab:threshold} and Figure~\ref{fig:threshold} report the results. The closed form describes the
full model well. Without a scale effect the two agree to within a few hundredths of a percentage point.
With one, the simulated threshold lies above the closed form when remediation is low, by 4 per cent at
$\nu=0.5$ and 13 per cent at $\nu=1$ without remediation, and the gap closes as remediation rises: at
$\tau_r=0.2$ the two agree to within $0.02$ percentage points.
This is part (iii) of Proposition~\ref{prop:threshold} at work. Population first falls, to about
two-thirds of its initial level, before habitat recovers, and an economy whose clean research would
sustain the stationary population once reached can lose decoupling during the undershoot. Panel~(b)
shows the paths: with clean research 3 per cent above the threshold population recovers towards its
stationary level, while 3 per cent below it population turns down again after nearly recovering, and
15 per cent below it decline accelerates as clean progress slows. Remediation narrows the undershoot,
which is why the gap between the simulated and closed-form thresholds closes as it rises.

Panel~B of the table shows the interaction between fertility and clean technology. Starting from a lower
taste for children, the economy needs a larger habitat per capita to restore replacement and settles at a
smaller population: from $\alpha_0=0.30$, the stationary population is a sixth of the initial one. With a
scale effect, it then needs far more clean research, $2.5$ times as much at $\nu=0.5$ and six times as
much at $\nu=1$.

\subsection{Semi-endogenous clean knowledge}
\label{sub:semiendog}
The threshold in Proposition~\ref{prop:threshold} is a condition on policy because clean progress is
proportional to effort. Under the alternative of a semi-endogenous law,
$B_{t+1}-B_t=B_t^{1-\beta_c}(R_tH_t)^{\lambda_c}$, the counterpart of the iso-elastic ideas law, it becomes
a condition on technology.

\begin{proposition}[Semi-endogenous clean knowledge]
\label{prop:semiendog}
Let $k\equiv\lambda_c/\beta_c$. A stationary population with growing output exists if and only if
$k\theta\ge1$, whatever the research share. If $k\theta<1$, population declines at
\[
\ln\tilde n=-\frac{\mu\,\zeta^\ast\ln(1+\bar g_A)\,(1-k\theta)}{1+\mu(1-k)} .
\]
\end{proposition}

\begin{proof}
Along a balanced path with constant $R$, $1+g_B=[\tilde n(1+g_h)]^{k}$ and, at the saturated limit,
$1+g_h=(1+\bar g_A)^{\sigma\theta/(1-\theta)}$. Pressure stays bounded at $\tilde n=1$ if and only if
$k\sigma\theta/(1-\theta)\ge\zeta^\ast$, that is $k\theta\ge1$. Otherwise, setting
$[(1+\bar g_A)^{\zeta^\ast}\tilde n/(1+g_B)]^{-\mu}=\tilde n$ and solving gives the rate.
\end{proof}

Under this law no permanent policy moves the threshold, since $R$ changes only the level of clean
knowledge, and with the returns used for general ideas, $k=\lambda/\beta=0.375$, the condition
$k\ge1/\theta=1.20$ fails. Decline is also faster than if clean progress did not depend on the size of
the research workforce, for which the denominator would be $1+\mu$. Whether a stationary population is a policy outcome thus
depends on whether clean innovation is fully or semi-endogenous, a question on which the evidence on
green innovation, rather than on fertility, will have to decide.

\subsection{Crowding in tastes, quality in costs}
\label{sub:split}
So far one index, habitat per capita $\epsilon=Q/D$, moves both the taste for children and their cost.
The evidence of Section~\ref{sub:pollution} suggests a finer split: stated ideals respond to crowding (Table~\ref{tab:prefs}) but show no detectable response to the
pollution load (Table~\ref{tab:pollution}), which points to quality acting mainly through
the costs of children. We nest both cases by letting quality enter tastes with weight $\iota\in[0,1]$,
\begin{align}
\alpha_t=\alpha\!\left(\frac{Q_t^{\iota}}{D_t}\right),
\qquad
\Phi_t=\Phi\!\left(\frac{Q_t}{D_t}\right),
\label{eq:split}
\end{align}
where $\Phi$ is the full cost of a child of Corollary~\ref{cor:goods}. The common index of
Section~\ref{sec:model} is $\iota=1$; $\iota=0$ lets crowding alone move tastes. Write
$\xi\equiv\mathrm{d}\ln\alpha/\mathrm{d}\ln(Q^{\iota}/D)$ and $\chi\equiv-\mathrm{d}\ln\Phi/\mathrm{d}\ln(Q/D)$, as
before.

\begin{proposition}[Split habitat]
\label{prop:split}
Let quality settle at $Q_\infty(\tau_r)$, with land area $\bar E$ fixed. Then:
\begin{enumerate}
\item[(i)] a stationary population $L^\ast$ solves
$\alpha\big(Q_\infty^{\iota}\bar E/L^\ast\big)=\big[\Gamma+\Phi(Q_\infty\bar E/L^\ast)\big]/
\big[1-\Phi(Q_\infty\bar E/L^\ast)\big]$; it exists if the difference of the two sides changes sign, and it
is then unique;
\item[(ii)] the feedback strength is $S^\ast\equiv-\mathrm{d}\ln n/\mathrm{d}\ln L$ at $L^\ast$,
$S^\ast=\mathcal{K}(\alpha^\ast)\xi+\chi$, and $L^\ast$ is locally stable if and only if $0<S^\ast<2$;
\item[(iii)] a cleaner habitat raises the stationary population with elasticity
\begin{align}
\frac{\mathrm{d}\ln L^\ast}{\mathrm{d}\ln Q_\infty}=\frac{\iota\,\mathcal{K}(\alpha^\ast)\xi+\chi}{S^\ast},
\label{eq:dLdQ}
\end{align}
which is one with the common index and $\chi/S^\ast$ when quality acts on costs only;
\item[(iv)] Proposition~\ref{prop:threshold}(i) and~(iii) and Corollary~\ref{cor:threshold}(c) hold with this
$L^\ast(\tau_r)$; the threshold is $R^\ast(\tau_r)=(g_B^\ast/\delta_B)L^\ast(\tau_r)^{-\nu}$, rising as
$L^\ast$ falls, and part (a) of the corollary becomes
\begin{align*}
\frac{\partial\ln R^\ast}{\partial\tau_r}=-\nu\,\frac{\iota\,\mathcal{K}(\alpha^\ast)\xi+\chi}{S^\ast}\,
\frac{Q_\infty'(\tau_r)}{Q_\infty(\tau_r)}.
\end{align*}
\end{enumerate}
\end{proposition}

\begin{proof}
By \eqref{eq:sol_n_hab} and \eqref{eq:S_def}, $\mathrm{d}\ln n=\mathcal{K}(\alpha)\xi\,
\mathrm{d}\ln(Q^{\iota}/D)+\chi\,\mathrm{d}\ln(Q/D)=(\iota\mathcal{K}\xi+\chi)\,\mathrm{d}\ln Q-S\,\mathrm{d}\ln L$.
Fertility therefore falls in $L$ at given quality, which gives (i), and the population map
$L_{t+1}=L_tn(L_t)$ has slope $1-S^\ast$ at $L^\ast$, which gives (ii). Setting $\mathrm{d}\ln n=0$ gives (iii).
The proofs of Proposition~\ref{prop:threshold}(i) and~(iii) use only that population converges to a
stationary $L^\ast$ once pressure vanishes, and that of Corollary~\ref{cor:threshold} only that
$\delta_BR(L^\ast)^{\nu}=g_B^\ast$; (iv) follows with (iii).
\end{proof}

The split changes one thing that matters. With the common index, a cleaner environment is converted
one for one into population; when quality acts on costs only, it is converted in proportion
$\chi/S^\ast$, about two-fifths at our estimates, and the rest raises habitat per capita and fertility
along the way. Remediation is then a weaker route to a larger population and a weaker substitute for clean
research. The calibration, however, offsets most of the difference for everything else. The pollution
coefficient of Section~\ref{sub:pollution} measures the response of fertility to the pollution load
relative to its response to density, which by the calculation in the proof identifies
$e_Q(\iota\mathcal{K}\xi+\chi)/S$ rather than $e_Q$ alone. Calibrated to that reduced form, a smaller
$\iota$ implies a larger $e_Q$, and the effect of ending pollution on the stationary population is
nearly unchanged. Section~\ref{sec:quant} reports both cases. The asymptotic rate of decline in
Proposition~\ref{prop:threshold}(ii) depends on the functional forms when $\iota<1$, and we do not
characterise it.

\section{Empirical evidence on the habitat--fertility link}
\label{sec:empirics}

The analysis above rests on the response of fertility to habitat per capita, through the taste for
children, $\alpha'(\epsilon)>0$, and through the time cost of rearing them, $\phi'(\epsilon)<0$. The existence of a stationary $\epsilon^\ast$ with replacement fertility, the stability
corridor $0<\epsilon^\ast n'(\epsilon^\ast)<2$ and the sign of the remediation-tax comparative
statics in \eqref{eq:dn_dtau} all follow from that response and its strength, and so, through the
stationary population, does the clean-research threshold \eqref{eq:R_star}; by
Proposition~\ref{prop:two_channels} and Corollary~\ref{cor:goods}, a response of either the taste or the time or goods cost of children suffices for stabilisation. This section sets out what can and cannot be learned about
that response, and about the parts of it due to tastes and to costs, from survey and national data.

The evidence supports four statements, and we organise the section around
them. All are statements about associations, not causal effects. First, fertility is positively
associated with proxies for habitat per capita. The gradient is monotone and precisely estimated in the
World Values Survey, reappears in the European Values Study, in two US surveys and in a panel of 161
countries, and, under the maintained conversion from settlement size to habitat per capita, implies a
demographic feedback inside the stable region of Proposition~\ref{prop:two_channels} (Sections~\ref{sub:gradient} and~\ref{sub:xi}). Second, the
gradient is not the work of a time cost alone. A pure time-cost channel leaves labour supply
unchanged, whereas household hours fall with habitat in the larger of our two US samples, though
not in the smaller (Section~\ref{sub:budget}). Third, preferences account for part of the gradient. The budget-based
test the model provides for this purpose leans on the homotheticity of logarithmic preferences,
which household budget data reject, and a version that relaxes it agrees with the taste channel in
sign only. Evidence outside the budget constraint is more informative: stated ideal fertility
rises with habitat at a little over half the rate of realised fertility; in the United States,
ideals respond to the habitat in which respondents grew up rather than the one in which they live,
although this runs through the family of origin; and in the main specification national and local
habitat move ideals at similar rates (Section~\ref{sub:prefs}). Under the interpretation of stated ideals as a taste measure, the point estimate attributes about
54 per cent of the gradient to tastes and the remainder to costs of children. Fourth, at a given density and income, fertility is higher where the pollution load on the habitat is lower, although stated ideals show no detectable response to pollution and reverse causality from fertility
to emissions cannot be excluded (Section~\ref{sub:pollution}). Section~\ref{sub:budget} summarises the
budget-based tests, whose derivation and estimates are in Appendix~\ref{app:budget}, and
Section~\ref{sub:caveats} states what remains unsettled.

%% -------------------------------------------------------------
\subsection{Data}
\label{sub:data}
%% -------------------------------------------------------------

Our source is the World Values Survey trend file 1981--2022, which pools seven waves and
$442{,}473$ respondents.\footnote{See \url{https://www.worldvaluessurvey.org/WVSEVStrend.jsp}.
Wave coverage is documented in the Common EVS/WVS Dictionary, which we use below; several items
are not fielded in every wave.} Habitat is proxied by \texttt{X049}, an eight-category measure of
the size of the respondent's town with thresholds at 2{,}000, 5{,}000, 10{,}000, 20{,}000,
50{,}000, 100{,}000 and 500{,}000 inhabitants. Desired fertility is \texttt{D017}, the ideal number
of children; realised fertility is \texttt{X011}, children ever had, top-coded at five. Controls
are age, sex, marital status, religiosity and, with the reservation noted below, education
and income.

Wave coverage is the binding constraint, and it differs sharply between the two fertility measures.
\texttt{D017} is fielded in WVS waves 1--4 only, and \texttt{X049} in waves 2--7, so the two
overlap in waves 2, 3 and 4: $98{,}512$ respondents in 51 countries and 77 country--waves, spanning
1989--2004. \texttt{X011}, by contrast, appears in all seven waves, giving $316{,}751$ respondents
in 100 countries and 231 country--waves once \texttt{X049} is also required. The ideal-fertility
measure is therefore confined to a historical window while realised fertility is not, which is one
reason we build the structural estimate of Section~\ref{sub:xi} on \texttt{X011}.

Three further features constrain the design. First, \texttt{X049} is suppressed for data protection
in EVS 2008 and EVS 2017, so those waves cannot contribute if the European file is appended; they
carry no \texttt{D017} in any case. Second, the archive documentation warns of cross-survey
inconsistencies in \texttt{X049}, and category boundaries differ in several country--waves; the
natural coarsening is \texttt{X049a}, a five-category version available in waves 2--7, rather than
\texttt{X050}, which is fielded only in waves 1 and 2 and cannot serve as a harmonised alternative.
Third, \texttt{X011} is top-coded at five children, which affects $8.1$ per cent of respondents and
compresses measured fertility in high-fertility cells.

The WVS carries neither education expenditure nor hours worked, so it cannot support the tests of
Appendix~\ref{sub:ltest}; what it can deliver is the gradient of Section~\ref{sub:gradient} and the
structural estimate of Section~\ref{sub:xi}, and those are what we report here. The test
on labour supply uses the US General Social Survey cumulative file, 1972--2024, which records
children ever born, the ideal number of children, size of place of residence and hours worked last
week; sampling-design variables are available from 1975 and size of place to 2022, which gives
$37{,}678$ respondents aged 25--54. The expenditure-share test requires budget data, for which we use the Interview Survey of
the US Consumer Expenditure Survey, data releases 2000--2025.

%% -------------------------------------------------------------
\subsection{Motivating fact: the settlement-size gradient}
\label{sub:gradient}
%% -------------------------------------------------------------

We first document the gradient that motivates the mechanism. For respondent $i$ in country $c$ and
wave $w$ we estimate
\begin{align}
y_{icw}
=\sum_{k=2}^{K}b_k\,\mathbf{1}\{\text{settlement bin}_{icw}=k\}
+\mathbf{x}_{icw}'\boldsymbol{\delta}
+\varsigma_{cw}+u_{icw},
\label{eq:gradient_spec}
\end{align}
where $y$ is \texttt{D017} or \texttt{X011}, $\mathbf{x}$ collects the controls listed above, and
$\varsigma_{cw}$ are country$\times$wave fixed effects, which remove any national shock common to all
settlement types. Observations are survey-weighted and standard errors are
clustered at the country--wave level. For realised fertility we restrict the sample to respondents
aged 45--70, for whom fertility is essentially complete. Figure~\ref{fig:gradient} plots
$\{\hat b_k\}$ against bin.

\begin{figure}[h!]\centering
\includegraphics[width=\textwidth]{fig_gradient.pdf}
\caption{Fertility by settlement size. Coefficients $\hat b_k$ from \eqref{eq:gradient_spec}
with the smallest settlement ($<2{,}000$) as the omitted category and 95\% confidence intervals,
country$\times$wave fixed effects, controlling for age, sex, marital status and religiosity.
Panel (a) is the ideal number of children (\texttt{D017}, WVS waves 2--4); panel (b) is children
ever had (\texttt{X011}) for respondents aged 45--70. Survey-weighted; standard errors clustered
at the country--wave level.}
\label{fig:gradient}
\end{figure}

The gradient is monotone across all eight bins and in both measures. Relative to settlements under
2{,}000 inhabitants, the ideal number of children falls by $0.053$ in towns of 2{,}000--5{,}000 and by
$0.330$ (s.e.\ $0.050$) in cities above half a million; completed fertility falls by $0.504$
(s.e.\ $0.051$) over the same range. This is the micro analogue of the comparative static
$\partial n/\partial \epsilon>0$ of Section~\ref{sec:model}.

The bad-control prediction of Appendix~\ref{sub:ltest} is visible in the same estimates. Adding
education and income, both chosen jointly with fertility given $\alpha$, attenuates the
largest coefficient from $-0.330$ to $-0.275$ for ideal fertility and from $-0.504$ to $-0.393$ for
realised fertility, a compression of roughly a fifth in each case. We therefore report the
specification without them as the headline and the attenuation alongside it rather than in its
place.

We stress what none of this establishes. Because every mechanism in Table~\ref{tab:fingerprints} generates the same pattern,
Figure~\ref{fig:gradient} is a motivating fact and not a test. Nor are settlement bins a measure of
$\epsilon_t=E_t/L_t$: they order respondents by the population of their locality, not by habitat per
person, and they do so in the cross-section rather than along the aggregate time path the model
describes. Section~\ref{sub:xi} converts settlement population into habitat per capita under an
assumed scaling of land area with population, and Appendix~\ref{sub:ltest} asks which mechanism could
generate the gradient. Neither addresses the aggregate time path, to which
Appendix~\ref{app:crosscountry} returns.

%% -------------------------------------------------------------
\subsection{Estimating $\xi$ by inverting the fertility rule}
\label{sub:xi}
%% -------------------------------------------------------------

The illustrative calibration of Section~\ref{sub:illustration} set the elasticity $\xi$ in
$\alpha(\epsilon)=\alpha_0\epsilon^{\xi}$ by an illustrative target: the value at which a doubling of
habitat per capita brings fertility to replacement, giving $\xi=0.146$. The data allow $\xi$ to be
estimated rather than imposed. The key observation is that the fertility rule
\eqref{eq:sol_n_hab} inverts in closed form.

\begin{proposition}[Recovering the fertility taste from observed fertility]
\label{prop:inversion}
For any $n<1/\phi$, the fertility rule $n=(\alpha-\Gamma)/[\phi(1+\alpha)]$ inverts to
\begin{align}
\alpha=\mathcal{A}(n)\equiv\frac{\phi n+\Gamma}{1-\phi n},
\qquad
\mathcal{A}'(n)=\frac{\phi(1+\Gamma)}{(1-\phi n)^{2}}>0 .
\label{eq:inversion}
\end{align}
At $n=1$, $\mathcal{A}(1)=(\phi+\Gamma)/(1-\phi)=\alpha^{\ast}$, the replacement-fertility taste of
Proposition~\ref{prop:two_channels}.
\end{proposition}

Given $(\phi,\Gamma)$ from the calibration, \eqref{eq:inversion} maps any observed fertility rate
into the taste parameter that rationalises it. Applying it cell by cell (we use country
$\times$ wave $\times$ settlement-bin cells, so that $n$ is a cell mean rather than an individual
count) yields $\hat\alpha_{j}$ for each cell $j$. The elasticity is then the slope of
\begin{align}
\ln \hat\alpha_{j}=\ln\alpha_0+\xi \ln \epsilon_{j}+\varsigma_{cw}+v_{j}.
\label{eq:xi_spec}
\end{align}

We implement \eqref{eq:xi_spec} on country$\times$wave$\times$bin cells built from \texttt{X011}
for respondents aged 45--70, taking $n$ as the cell's weighted mean number of children divided by two\footnote{Dividing by
$2.1$ instead, as with the UN data below, gives a composite of $0.0178$ rather than $0.0185$.} so that
$n=1$ is replacement, and keeping cells with at least thirty respondents. This yields $1{,}031$
cells across $225$ country--waves. The mean implied taste, weighting cells by size, is $\hat\alpha=0.493$, above both
$\alpha_0$ and $\alpha^{\ast}$, as one expects in a global sample whose older cohorts completed
their fertility above replacement.

Implementing \eqref{eq:xi_spec} requires $\epsilon$ in levels, which the settlement bins do not
provide: they record the \emph{population} of a locality, not space per person. Writing
$\varrho$ for the elasticity of a locality's land area with respect to its population, density scales
as population to the power $1-\varrho$, and the regression of $\ln\hat\alpha_j$ on minus the log of
bin-midpoint population identifies $\xi(1-\varrho)$ rather than $\xi$ itself. That composite is
estimated precisely,
\begin{align}
\widehat{\xi(1-\varrho)}=0.0185\quad(\text{s.e.}\ 0.0012),
\label{eq:xi_composite}
\end{align}
and is stable across the choices the construction forces on us: it lies between $0.017$ and $0.022$
when the open top bin is assigned a midpoint of $0.75$, $1$ or $2$ million, when the minimum cell
size is $15$ or $50$, when the age window is $40$--$75$ or $50$--$70$, when the sample is restricted
to women, and when it is split into waves 2--4 and waves 5--7. The $\hat\alpha_j$ are also monotone
in bin, falling by $0.129$ log points from the smallest settlements to cities above half a million.

Recovering $\xi$ itself requires a value for $\varrho$. Land area scales sublinearly with population,
and we take $\varrho=0.85$ as a benchmark, reporting $0.80$ and $0.90$ alongside it.\footnote{The
conversion is an assumption rather than an estimate.} At $\varrho=0.85$, \eqref{eq:xi_composite}
implies
\begin{align}
\hat\xi=0.123\qquad[0.108,\,0.139],
\label{eq:xi_estimate}
\end{align}
against the $0.146$ imposed in Section~\ref{sub:illustration}; the calibrated value is reproduced
exactly at $\varrho=0.873$, and the estimate is $0.093$ at $\varrho=0.80$ and $0.185$ at
$\varrho=0.90$.

When the time cost also responds to habitat, the inversion holds $\phi$ fixed and therefore attributes the whole fertility
gradient to tastes. Since $\mathrm{d}\ln\mathcal{A}(n)/\mathrm{d}\ln n=1/\mathcal{K}$ at the
stationary point, what that procedure recovers is
\begin{align}
\hat\xi\;\longrightarrow\;\xi+\frac{\chi}{\mathcal{K}(\alpha^{\ast})},
\label{eq:xi_bound}
\end{align}
the total habitat elasticity of fertility expressed in taste units. It is an upper bound on $\xi$,
tight only when the time-cost channel is absent, and it is exactly the quantity that governs
stability, because $S=\mathcal{K}\left(\xi+\chi/\mathcal{K}\right)$.

Two qualifications govern how this number should be read. The inference from
\eqref{eq:xi_composite} to \eqref{eq:xi_estimate} rests on $\varrho$ and not on anything the survey
measures. And, by \eqref{eq:xi_bound}, \eqref{eq:xi_estimate} estimates $\xi+\chi/\mathcal{K}$: the total habitat
elasticity of fertility in taste units, and an upper bound on $\xi$ itself. A time-cost elasticity
of $\chi=0.1$ would lower the implied taste elasticity to $0.070$. Separating the two requires the
expenditure-share regression of Appendix~\ref{sub:ltest} or the ideal-fertility ratio of
Section~\ref{sub:prefs}, through the decomposition \eqref{eq:decomposition}. Without taking a view on the mix, and under the maintained conversion from settlement population
to habitat per capita, the estimated association maps into an implied feedback strength $S$ of
Proposition~\ref{prop:two_channels}, which is simply the habitat elasticity of fertility.
Regressing $\ln n_j$ directly on the same regressor gives $\widehat{S(1-\varrho)}=0.0313$
(s.e.\ $0.0020$), hence at $\varrho=0.85$
\begin{align}
\hat S=0.209\qquad[0.182,\,0.236],
\label{eq:S_estimate}
\end{align}
which implies monotone convergence with a half-life of about three generations. The estimate moves
to $0.16$ at $\varrho=0.80$ and $0.31$ at $\varrho=0.90$, remaining far from the oscillation threshold
throughout. Read through \eqref{eq:S_def} at the sample mean of $\hat\alpha$, where
$\mathcal{K}=1.57$, it bounds the taste elasticity at $0.133$ if $\chi=0$, and implies $0.069$ at
$\chi=0.1$ and zero at $\chi=0.21$. These figures differ from $0.123$ and $0.070$ above because the
inversion is applied cell by cell, so that the two composite estimates weight $\mathcal{K}$ across
cells rather than evaluating it at a single point. Their ratio, $0.0313/0.0185$, implies an effective
$\mathcal{K}$ of $1.69$, between $1.57$ at the sample mean and $1.88$ at $\alpha^\ast$.

\paragraph{A country--year panel.}
The UN World Population Prospects 2024 provide a check that needs no assumption about $\varrho$,
because national population density is observed directly. For 161 countries with at least one
million inhabitants over 1950--2023 we apply the inversion to $n=\text{TFR}/2.1$ and regress
$\ln\hat\alpha$ on $\ln(1/\text{density})$. The slope is $0.059$ (s.e.\ $0.017$) in the 2023
cross-section and $0.103$ (s.e.\ $0.019$) in the pooled sample with year effects; with country and
year effects it is $0.243$ (s.e.\ $0.035$), rising to $0.381$ (s.e.\ $0.032$) after 1990. Every
estimate has the predicted sign. The between-country estimates lie below the survey-based
$0.123$ and the within-country estimates above it, as one would expect: the former are attenuated
by the crudeness of national density as a measure of habitat, while the latter load the
demographic transition, during which density rose and fertility fell for reasons held fixed in the
model, onto the habitat elasticity. We regard the survey estimate, bracketed by the two, as the
more credible.

Had the estimated response come out substantially smaller, the consequence would have been more than
quantitative: by Proposition~\ref{prop:two_channels}, local stability of $\epsilon^{\ast}$ requires
$0<\epsilon^{\ast}n'(\epsilon^{\ast})<2$, and a weak habitat response places the economy near the
lower edge of that corridor, so that the self-correcting demography on which our main result depends
would operate only slowly. At $\hat S\approx0.21$ the estimated total response is close to the
illustrative $S=0.275$ of Section~\ref{sub:illustration}, comfortably inside the corridor and well short of the oscillation threshold $S=1$ of
Proposition~\ref{prop:two_channels}, whichever channel accounts for it.

%% -------------------------------------------------------------
\subsection{Tastes or costs: what household budgets show}
\label{sub:budget}
%% -------------------------------------------------------------

A negative association between settlement size and fertility is implied by essentially every model in
the literature we build on, most of which act through the cost of children \citep{DAO2013442,dao25}. The
gradient of Section~\ref{sub:gradient} is therefore consistent with our mechanism but no evidence for it
over these alternatives. The closed-form household solution separates the channels in principle. Under
logarithmic preferences the consumption share, total education time and the education spending share
depend on the taste for children alone, and every cost of children, in time, in goods or through the
price of education, leaves them unchanged (Proposition~\ref{prop:identification} in
Appendix~\ref{app:budget}). The ratio of their habitat elasticity to that of fertility is
$-(\alpha-\Gamma)/(1+\Gamma)\in[-0.160,-0.126]$ when habitat acts through tastes and zero when it acts
through costs. Labour supply discriminates less finely: a pure time cost leaves it unchanged, while a
goods cost such as housing moves it in the same direction as a taste shift.

Appendix~\ref{app:budget} takes these restrictions to the US Consumer Expenditure Survey and the General
Social Survey. Both surveys reproduce the fertility gradient, with feedback strengths of $0.25$ and
$0.16$ against $0.21$ in the World Values Survey. Household hours fall with habitat in the larger survey,
with ratios between $-0.29$ and $-0.39$ whose intervals exclude zero, which rejects a pure time-cost
channel within the maintained labour-supply specification; the smaller survey is not powered to detect
a response either way. The share restriction fares worse. It relies on the unit income elasticities of
logarithmic preferences, which budget data reject: consumption relative to the wage falls with the wage,
and the estimated ratio changes sign with the treatment of the wage. A Stone--Geary version that relaxes
homotheticity agrees with the taste channel in sign, but its estimates are an order of magnitude too
large to be attributed to tastes alone. The budget evidence thus rules out a pure time-cost account and
is suggestive, no more, of a taste component. The rest of the section turns to evidence outside the
budget constraint.

%% -------------------------------------------------------------
\subsection{Direct evidence on preferences}
\label{sub:prefs}
%% -------------------------------------------------------------

The budget tests of Section~\ref{sub:budget} work through budget constraints, and they reach their limit
where a goods cost and a taste shift move the budget in the same direction. A different route
exploits the fact that preferences and costs leave different traces outside the budget. A cost of
children depends on where a household lives now and binds on what it does; a taste is carried by
the person, may be formed earlier, and shows up in what people say they want as well as in what
they do. We use three such traces.

\paragraph{Ideal and realised fertility.}
The World Values Survey asks the same respondents how many children they have and how many they
would ideally want. The model contains no stated-ideal variable, so reading the two together requires
an additional assumption. We assume that the stated ideal is the fertility the respondent would choose
given her taste $\alpha(\epsilon)$ with the costs of children held at a common reference level,
$n^{\text{ideal}}=n\big(\alpha(\epsilon);\bar\Phi\big)$, so that it obeys the same taste-to-fertility
mapping as realised fertility but does not respond to costs. Under this assumption, if habitat acts
through tastes alone both measures rise with it at the same proportional rate; if it acts through costs
alone the stated ideal does not move and the shortfall of realised below ideal fertility widens as
habitat shrinks. In the two-channel model of Section~\ref{sec:model} the ratio of the two habitat
elasticities is then the taste share of the fertility gradient,
\begin{align}
\varpi\;\equiv\;\frac{\partial\ln n^{\text{ideal}}/\partial\ln\epsilon}{\partial\ln n/\partial\ln\epsilon}
=\frac{\mathcal{K}\xi}{\mathcal{K}\xi+\chi},
\label{eq:tau}
\end{align}
the same object that $\Lambda'$ was designed to recover in \eqref{eq:Lambda_share}, now measured
without any assumption on the form of preferences over consumption but conditional on the
interpretation of stated ideals.

We estimate both elasticities on respondents aged 45--70, for whom fertility is complete, applying
the same top-code of five to both measures. To the World Values Survey (waves 2--4) we add the
European Values Study trend file, whose 1990 wave asks about ideal fertility in 29 European countries.
Of these, 23 enter the sample: we drop two surveys that the World Values Survey already contains
(Czechoslovakia and Spain), three that lack size of town (Finland, Iceland and Lithuania) and one
that lacks religiosity (Estonia). Sweden's 1999 survey, the only later one that asks about ideal
fertility, adds a twenty-fourth; its interviews enter through the European file only, since the World
Values Survey file contributes none of them. The pooled sample covers 99 country--waves, 75 from the World Values
Survey and 24 from the European survey; the 75 are the 77 of Section~\ref{sub:data} less China 1995 and
Pakistan 1997, which do not record religious attendance. Panel~A of
Table~\ref{tab:prefs} reports the result. Ideal fertility rises with habitat at a little over half
the rate of realised fertility: $\hat\varpi=0.54$, with an interval of $[0.43,\,0.68]$ that excludes
both a pure cost account, $\varpi=0$, and a pure taste account, $\varpi=1$. The World Values Survey
alone gives $0.59$ $[0.46,\,0.75]$ and the European sample alone $0.38$ $[0.18,\,0.75]$. The gap between ideal and realised fertility correspondingly widens with density. An
obvious concern is that respondents who have completed their families report as ideal the number
they have. The ideal-fertility gradient is, however, almost the same among respondents aged
18--34, most of whom have not, as among those aged 45--70, which is not what ex-post
rationalisation would produce.

Combined with the feedback strength $\hat S=0.209$ of \eqref{eq:S_estimate}, $\hat\varpi$ delivers
the decomposition \eqref{eq:decomposition} directly. With $\mathcal{K}=1.57$ at the sample mean
of $\hat\alpha$,
\begin{align}
\hat\xi=\frac{\hat\varpi\hat S}{\mathcal{K}}=0.072\quad[0.053,\,0.091],
\qquad
\hat\chi_c=(1-\hat\varpi)\hat S=0.096\quad[0.067,\,0.125],
\label{eq:decomp_estimate}
\end{align}
where the intervals combine the sampling uncertainty in $\hat\varpi$ and in $\hat S$ by the delta
method, treating the two estimates as independent. Since both draw on the World Values Survey, a joint
bootstrap would be preferable; it is not straightforward here because $\hat S$ is estimated on cells
and $\hat\varpi$ on individuals. The cost elasticity $\hat\chi_c$ collects every cost of children, time
and goods alike, and corresponds to $\chi_\Phi$ in Corollary~\ref{cor:goods}. Under the interpretation
of stated ideals as a taste measure, the point estimate attributes about 54 per cent of the observed
gradient to tastes and the remainder to costs, and their combined strength lies far below the oscillation threshold of
Proposition~\ref{prop:two_channels}. If stated ideals partly internalise the costs respondents
expect to face, $\hat\varpi$ overstates the taste share and \eqref{eq:decomp_estimate} is an upper
bound for $\xi$; if ideals are measured with more error than realised fertility, it understates
it.

\paragraph{Habitat in childhood.}
A cost of children is paid where children are raised, whereas a taste may be formed where the
parent grew up. The General Social Survey records the type of place in which respondents lived at
age 16, alongside the size of place in which they live now. Panel~B of Table~\ref{tab:prefs}
reports both elasticities for completed fertility at ages 40--64, with habitat at 16 measured as
minus the log of a representative population for each type of place. Childhood habitat predicts
fertility conditional on current habitat, and the comparison of movers with stayers is telling:
among respondents who have left the city in which they grew up, and for whom the two habitats are
therefore distinct, the childhood elasticity of $0.012$ is as large as the current one. Stated
ideals make the point more sharply still. The ideal number of children does not respond to current
size of place at all, but it does respond to the place in which the respondent grew up.

These patterns are hard to reconcile with costs, which depend on where a household lives and not
on where its adults were raised. They have a qualification of their own, however. Both childhood
effects vanish or reverse once the respondent's number of siblings and parents' education are held
fixed, and the coefficient for completed fertility turns slightly negative. Habitat in childhood therefore acts through the family of origin: respondents raised in
roomier places come from larger families, and the preference for a larger family is transmitted.
Within the model this is what one would expect if the parents' taste was itself shaped by their
habitat, but the data cannot distinguish that account from one in which family size is
transmitted for reasons unrelated to habitat.

\paragraph{National habitat per capita.}
The model's habitat is economy-wide, $\epsilon_t=E_t/L_t$, whereas the survey measures are local.
Merging the World Values Survey and the European Values Study with national population density from
the UN World Population Prospects and with GDP per capita from the World Bank lets us include both.
The pooled sample has ideal fertility for 79 countries, eleven of them observed only in the
European survey. Panel~C of Table~\ref{tab:prefs} reports the result. Across countries, and
controlling for income, national habitat raises the ideal number of children with an elasticity of
$0.039$ (s.e.\ $0.011$). The estimate is unchanged when continent effects are added, which removes
the broad regional contrasts in culture and religion that a cross-country comparison would
otherwise load onto density, and it is of the same size, $0.046$ (s.e.\ $0.016$), among the 37
European countries alone. When local habitat is included the two enter together, at $0.051$
nationally and $0.0074$ locally. Converting the local coefficient from settlement population to
habitat per capita at $\varrho=0.85$ gives $0.049$, so that national and local habitat move ideals
at essentially the same rate, the condition under which the cross-sectional evidence speaks to
the aggregate object in the model, and which Section~\ref{sub:caveats} would otherwise have to
leave untested. With both continent effects and local habitat the national coefficient is
$0.045$ (s.e.\ $0.024$), borderline at conventional levels. Within countries, the change in
density over the survey window is too small to identify anything.

\begin{table}[h!]\centering
\begin{threeparttable}
\caption{Direct evidence on the taste for children}
\label{tab:prefs}
\footnotesize\setlength{\tabcolsep}{3.5pt}
\begin{tabular}{lrcc}
\toprule
\multicolumn{4}{l}{\emph{A. Ideal and realised fertility, World Values Survey and European Values Study}}\\
& $N$ & Elasticity & \\
\midrule
Ideal, ages 45--70, pooled                & $40{,}331$ & $0.0148$ $(0.0022)$ & \\
Realised, ages 45--70, pooled             & $40{,}331$ & $0.0274$ $(0.0026)$ & \\
Ratio $\hat\varpi$, pooled                  &            & $0.539$ & $[0.428,\,0.678]$\\
Ratio $\hat\varpi$, WVS only                & $29{,}715$ & $0.586$ & $[0.463,\,0.754]$\\
Ratio $\hat\varpi$, EVS only                & $10{,}616$ & $0.377$ & $[0.181,\,0.751]$\\
Gap, ideal minus realised (WVS, children) & $29{,}715$ & $-0.0260$ $(0.0059)$ & \\
Ideal, ages 18--34 (WVS)                  & $39{,}617$ & $0.0150$ $(0.0017)$ & \\
Ideal, ages 35--44 (WVS)                  & $21{,}120$ & $0.0175$ $(0.0029)$ & \\
\midrule
\multicolumn{4}{l}{\emph{B. Habitat now and at age 16, General Social Survey}}\\
& $N$ & Current & At age 16\\
\midrule
Completed fertility, 40--64               & $25{,}790$ & $0.0123$ $(0.0031)$ & $\phantom{-}0.0089$ $(0.0022)$\\
\quad movers                             & $16{,}413$ & $0.0132$ $(0.0037)$ & $\phantom{-}0.0120$ $(0.0025)$\\
\quad stayers                            & $8{,}994$  & $0.0113$ $(0.0047)$ & $\phantom{-}0.0052$ $(0.0044)$\\
\quad + siblings, parents' education     & $18{,}104$ & $0.0125$ $(0.0037)$ & $-0.0051$ $(0.0027)$\\
Ideal number, 25--64                     & $18{,}636$ & $0.0003$ $(0.0017)$ & $\phantom{-}0.0043$ $(0.0014)$\\
\quad + siblings, parents' education     & $18{,}636$ & $0.0009$ $(0.0017)$ & $-0.0004$ $(0.0014)$\\
\midrule
\multicolumn{4}{l}{\emph{C. National and local habitat, ideal fertility, WVS and EVS with UN WPP and World Bank}}\\
& $N$ & National & Local\\
\midrule
Wave effects, 79 countries                & $185{,}255$ & $0.0386$ $(0.0109)$ & \\
\quad + continent effects                 & $185{,}255$ & $0.0383$ $(0.0165)$ & \\
\quad Europe only, 37 countries           & $81{,}692$  & $0.0456$ $(0.0157)$ & \\
With local habitat, 61 countries          & $120{,}933$ & $0.0509$ $(0.0165)$ & $0.0074$ $(0.0040)$\\
\quad + continent effects                 & $120{,}933$ & $0.0454$ $(0.0241)$ & $0.0096$ $(0.0029)$\\
\quad Europe only, 34 countries           & $62{,}354$  & $0.0376$ $(0.0204)$ & $0.0134$ $(0.0022)$\\
\bottomrule
\end{tabular}
\begin{tablenotes}\footnotesize
\item Poisson pseudo-maximum-likelihood elasticities, survey-weighted, with controls for age, sex,
marital status and religious attendance. Panel~A: WVS waves 2--4 and EVS waves 1990 and 1999,
country$\times$wave effects, both measures top-coded at five, bootstrap over country--waves. Panel~B: year effects and the controls of
Table~\ref{tab:gss}; habitat is minus the log of size of place (current) or of a representative
population for the type of place at 16; movers left the city in which they lived at 16; the
ideal-fertility rows use respondents with siblings and parents' education recorded, so that both rows
share one sample; bootstrap over sampling clusters. Panel~C: national habitat is minus the log of population density in the
survey year (UN WPP 2024); local habitat is minus the log of settlement population; both
regressions control for log GDP per capita (World Bank) and include study$\times$wave effects;
bootstrap over countries.
\end{tablenotes}
\end{threeparttable}
\end{table}

Taken together, the three traces point the same way. Habitat moves what people say they want, not
only what they do; it does so whether the habitat is the one they live in now or the one they grew
up in; and, in the main specification, stated ideals respond to national and local habitat at
similar rates. None of these is
decisive on its own, and each rests on an interpretation: of stated ideals as preferences, of
childhood habitat as distinct from family background, of between-country variation as more than
regional culture. But a pure cost account has to explain all three away, and, under the interpretation of stated ideals as a taste measure, the decomposition in
\eqref{eq:decomp_estimate} attributes about 54 per cent of the habitat gradient to tastes.

%% -------------------------------------------------------------
\subsection{Crowding and pollution}
\label{sub:pollution}
%% -------------------------------------------------------------

The evidence so far measures one of the two channels through which habitat per capita falls. Settlement
size and national density record how many people share a habitat, that is $D_t$ in
\eqref{eq:eps_def}, but not how polluted it is. The pollution channel requires a measure of quality
$Q_t$. With $\ln\epsilon=\ln Q-\ln D$ and quality falling with the pollution load per unit of land,
$\ln Q=-e_Q\ln(\text{load}/\bar E)$ with $e_Q>0$, the model gives
\begin{align}
\ln\alpha=\text{const}+\xi\,(-\ln D)+\xi e_Q\,\big(-\ln(\text{load}/\bar E)\big),
\label{eq:pollution_spec}
\end{align}
so that, holding density fixed, fertility and the implied taste for children should rise as the
pollution load falls. Written in terms of $\alpha$, \eqref{eq:pollution_spec} assumes the common index and
no cost response. The implied taste $\hat\alpha$ of Section~\ref{sub:xi} records the total response of
fertility in taste units, as in \eqref{eq:xi_bound}, and in general the ratio of the pollution
coefficient to the density coefficient identifies $e_Q(\iota\mathcal{K}\xi+\chi)/S$, which equals $e_Q$
under the common index (Section~\ref{sub:split}). We measure the load by fossil and industrial carbon-dioxide emissions per square
kilometre of land, from the Global Carbon Budget as compiled by Our World in Data, and merge it with
the UN World Population Prospects panel and with the World Values Survey. Carbon dioxide is not itself
a local pollutant, but it is emitted together with the particulates, sulphur and nitrogen oxides that
are, and it is the one emissions series available for all countries since 1950. Because emissions rise
with industrialisation, every specification controls for log GDP per capita.

\begin{table}[h!]\centering
\begin{threeparttable}
\caption{Crowding and pollution}
\label{tab:pollution}
\footnotesize
\begin{tabular}{lcccc}
\toprule
 & Crowding & Cleaner & & \\
 & $-\ln D$ & $-\ln(\text{CO}_2/\text{km}^2)$ & $N$ & Countries\\
\midrule
\multicolumn{5}{l}{\emph{A. Implied taste $\ln\hat\alpha$, UN World Population Prospects, 1950--2022}}\\
Pooled, year effects             & $-0.018$ $(0.021)$ & $0.087$ $(0.017)$ & 9,821 & 153\\
Country and year effects         & $\phantom{-}0.302$ $(0.044)$ & $0.011$ $(0.019)$ & 9,821 & 153\\
Country and year effects, 1990--2022 & $\phantom{-}0.289$ $(0.036)$ & $0.068$ $(0.016)$ & 4,947 & 153\\
\midrule
\multicolumn{5}{l}{\emph{B. World Values Survey, national measures}}\\
Completed fertility, ages 45--70           & $-0.066$ $(0.041)$ & $0.096$ $(0.041)$ & 131,391 & 99\\
\quad + local habitat                      & $-0.093$ $(0.038)$ & $0.139$ $(0.036)$ & 102,555 & 94\\
Ideal number of children                   & $\phantom{-}0.021$ $(0.041)$ & $0.027$ $(0.032)$ & 138,882 & 67\\
\quad + local habitat                      & $\phantom{-}0.071$ $(0.054)$ & $0.014$ $(0.037)$ & 92,198 & 49\\
\bottomrule
\end{tabular}
\begin{tablenotes}\footnotesize
\item All specifications control for log GDP per capita. Panel A: OLS of $\ln\mathcal{A}(\text{TFR}/2.1)$,
countries with at least one million inhabitants, standard errors clustered by country. Panel B: Poisson
pseudo-maximum likelihood with wave effects and the controls of Table~\ref{tab:prefs}, completed
fertility in waves 1--7 and ideal fertility in waves 1--4, bootstrap over countries. Because
$-\ln(\text{CO}_2/\text{km}^2)=-\ln(\text{CO}_2/L)-\ln D$, the crowding coefficient measures the effect of
density over and above its contribution to the pollution load. Emissions: Global Carbon Budget via
Our World in Data.
\end{tablenotes}
\end{threeparttable}
\end{table}

Table~\ref{tab:pollution} reports the result. A cleaner habitat goes with higher fertility at a given
density and income. Across countries the coefficient is $0.087$ (s.e.\ $0.017$), and in the World
Values Survey completed fertility rises with it at $0.096$, or $0.139$ when local habitat is also
controlled, so that the pollution load and settlement size enter together. Within countries the
association is absent over the whole period but appears after 1990, at $0.068$ (s.e.\ $0.016$), the
years in which the emission series is most reliable and environmental regulation most varied. The
crowding coefficient is the effect of density beyond its contribution to the load; across countries it
is absorbed by the load, while within countries it carries the demographic-transition component noted
in Section~\ref{sub:xi}.

Stated ideals tell a different story. Ideal fertility does not respond to the national pollution load,
with coefficients of $0.027$ and $0.014$, neither distinguishable from zero. It does respond to national
density when the load is not controlled (Table~\ref{tab:prefs}, Panel~C). With the load controlled, the
density coefficients for ideals, $0.021$ and $0.071$, are positive but imprecise. The absence of a detectable ideal-fertility response is
consistent with pollution operating mainly through costs rather than stated preferences: a polluted
habitat is associated with fewer children without fewer wanted, as one would expect if pollution acts through child health and the time and money needed to
raise children in a degraded environment \citep{mari10,Varvarigos2013}, the cost channel of
Section~\ref{sec:model} and Corollary~\ref{cor:goods}. For crowding, the evidence that ideals move as well as realised fertility comes from
Table~\ref{tab:prefs} and not from Table~\ref{tab:pollution}, where the crowding coefficient measures
only the effect of density beyond its contribution to the load. Population stabilisation does not depend on
which, by Proposition~\ref{prop:two_channels} and Corollary~\ref{cor:goods}; Section~\ref{sub:split} shows
what the split implies for the stationary population and for remediation. Three cautions apply. First, emission intensity at a given
income reflects industrial structure and energy sources, which differ across countries for reasons that
may also bear on fertility. Second, carbon dioxide measures pressure on the habitat, not its quality.
Third, and most important, is reverse causality. \citet{beulmann2025} document that within OECD countries since 1835 lower fertility has gone
together with higher carbon-dioxide emissions, and explain it by a causal path running from fertility to
emissions: fewer children receive more education, and the larger stock of human capital raises output and
emissions. That mechanism alone would generate the association in Table~\ref{tab:pollution} without any
effect of pollution on fertility. Two features of our estimates speak against it as the whole explanation.
The association with completed fertility holds at a given income, which absorbs much of the human-capital
channel, and it appears across countries and within countries after 1990; the absence of any response in
stated ideals, however, is also what the reverse channel would predict. We therefore read the pollution
estimates as consistent with, but not evidence for, the pollution channel. A sharper test would use
direct measures of local environmental quality, such as fine-particulate concentrations, matched to the
location of survey respondents, and exploit variation in pollution that is unrelated to local fertility.

%% -------------------------------------------------------------
\subsection{What these exercises cannot settle}
\label{sub:caveats}
%% -------------------------------------------------------------

Four limitations should be stated plainly.

First, the relationship is simultaneous by construction. This section attempts to
measure $\alpha(\epsilon)$, but the model of Section~\ref{sec:model} makes $\alpha_t$ the solution to a
within-period fixed point in which habitat depends on output and output depends on labour supply,
which depends on $\alpha_t$. The empirical counterpart of that simultaneity is residential sorting:
households that want children move to roomier places. Our estimates are therefore associations
consistent with the mechanism, not causal estimates of it, and we do not claim otherwise.

Second, $\epsilon$ in the model is an economy-wide aggregate, while settlement size is local.
Reading the cross-sectional gradient as evidence on $\alpha(\cdot)$ requires assuming that the same
preference function maps local and national roominess into the taste for children. Panel~C of
Table~\ref{tab:prefs} provides some support: national habitat raises stated ideals across 79
countries, with continent effects and within Europe, and in the main specification national and
local habitat move ideals at similar rates. The two diverge, however, when continent effects are
added or the sample is restricted to Europe, and the national estimate, identified from differences
between countries, becomes borderline once continent effects and local habitat are included
together. The pollution channel is measured more crudely still: the evidence of
Section~\ref{sub:pollution} rests on national carbon-dioxide emissions per square kilometre, a proxy for
the pollution load on the habitat rather than a measure of its quality, and it shows up in realised
fertility but not in stated ideals.

Third, the case for a taste component rests on evidence from outside the budget constraint, and
each piece carries an interpretation. The decomposition reads stated ideals as preferences
untouched by costs; to the extent that respondents report the family they expect to afford, it
overstates the taste share. The childhood-habitat evidence works through the family of origin and
cannot separate a taste shaped by the parents' habitat from family size transmitted for other
reasons. And the national-habitat evidence comes from differences between countries, which survive
continent effects and hold within Europe but cannot exclude every national difference correlated
with density. The test that would settle the question within the model, based
on expenditure shares, depends on a preference specification the data reject. We therefore regard
the taste channel as supported, with a point estimate of its share of about one half, but not as established beyond the
reach of a cost account that is willing to explain each of these patterns separately.

Fourth, the evidence is cross-sectional, while the model describes an aggregate time path. Over time
the habitat channel accounts for little of the fertility decline: Korea's population doubled between
1960 and 2023, which at the illustrative $\xi=0.146$ lowers the model's total fertility rate from $1.71$ to $1.33$, and by
less at the estimated elasticities of Section~\ref{sec:quant}, whereas
observed fertility fell from $6.0$ to $0.72$ (Appendix~\ref{app:crosscountry}). Most of the demographic
transition operates through forces held fixed in the model, and the central dynamic prediction, that
population decline raises habitat per capita and with it fertility, has not yet been observed on a
scale that would show up in the data.

% ============================================================
\section{Quantitative implications}
\label{sec:quant}
% ============================================================

We now calibrate the habitat block to the estimates of Section~\ref{sec:empirics} and ask what the
model implies for nine economies whose fertility is below replacement, in most of them far below: how much clean research each
needs for its population to settle, how that compares with the clean progress it achieves today, and at
what level its population would settle.

\subsection{Calibration}
\label{sub:recal}
Table~\ref{tab:recal} lists the parameters. The household and technology parameters are those of
Section~\ref{sub:illustration}. Four sets of parameters govern the habitat block and clean technology.

\emph{Feedback strength and its composition.} The feedback strength is set to the survey estimate
$\hat S=0.209$ and split between tastes and costs with the ideal-fertility share $\hat\varpi=0.54$
(Section~\ref{sub:prefs}). Writing $\alpha(\epsilon)=\alpha_i\epsilon^{\xi}$ and
$\phi(\epsilon)=\phi_0\epsilon^{-\chi}$, where $\phi$ stands for the full cost of a child
(Corollary~\ref{cor:goods}), the taste elasticity solves $\mathcal{K}(\alpha^\ast)\xi=\hat\varpi\hat S$
and the cost elasticity is $\chi=(1-\hat\varpi)\hat S$, which gives $\xi=0.060$ and $\chi=0.096$. We
evaluate $\mathcal{K}$ at the baseline stationary taste $\alpha^\ast=0.432$ rather than at the sample mean
of $\hat\alpha$ used in \eqref{eq:decomp_estimate}, because the feedback that governs stabilisation
operates at the stationary point; at the sample mean the taste elasticity would be $0.072$. Because the
cost is normalised at 2023 rather than at $\epsilon^\ast$, each country's own stationary taste lies
somewhat below $0.432$, and the feedback strength at its stationary point somewhat above $0.209$, from
about $0.22$ in France to $0.25$ in Korea, far inside the monotone region of
Proposition~\ref{prop:two_channels}. The taste elasticity is two-fifths of the illustrative $0.146$; the
total feedback is only somewhat below the illustrative $0.275$, because the cost channel makes up most of
the shortfall.

\emph{Pollution.} Environmental quality falls with the pollution load per square kilometre,
$Q=\big[(1+c_0)/(1+c_0x)\big]^{\mu}$, where $x$ is the load relative to its 2022 value. The local
elasticity at $x=1$, $e_Q=\mu c_0/(1+c_0)$, is set to the within-country estimate of
Section~\ref{sub:pollution}, $e_Q=0.068/0.289=0.235$, the ratio of the pollution and density
coefficients for 1990--2022. Because the within-country density coefficient also loads the demographic
transition (Section~\ref{sub:xi}), this ratio is if anything too low, which is one reason
Table~\ref{tab:sens} also reports $e_Q=1$. The curvature $c_0$ fixes the largest gain in quality that ending
pollution could deliver, $(1+c_0)^{\mu}$. Our data do not identify it; we set $c_0=4$, a maximum gain of
$61$ per cent, and report $c_0=1$ and $9$ below; Appendix~\ref{app:calib} explains the choice.
Together $e_Q$ and $c_0$ imply $\mu=e_Q(1+c_0)/c_0=0.29$, the elasticity of quality with respect to a
large load, against $\mu=1$ in the illustration of Section~\ref{sub:illustration}. This calibration uses the
common index, $\iota=1$. If quality acts on costs only, $\iota=0$, the same coefficient ratio identifies
$e_Q\chi/S$ (Proposition~\ref{prop:split}), so that $e_Q=0.235\times0.209/0.096=0.51$; we report that case
alongside the baseline.

\emph{Clean research.} Each country's current clean progress is read from its 2005--2022 trend in
carbon-dioxide emissions per unit of GDP, $g_{\text{int},i}$, which is negative in every country we
consider: per generation, $g_{B,i}=(1+g_{\text{int},i})^{-30}-1$. With population normalised to one in
2023, current clean progress is $\delta_BR_i=g_{B,i}$, so clean effort can be measured relative to the
effort that would produce the current measured progress, $m\equiv R/R_i$, without data on research
spending. Because the decline in emission intensity also reflects regulation, energy prices, structural
change and imported technology, $R_i$ is an effort equivalent, not a measure of research spending. The threshold \eqref{eq:R_star} then
becomes
\begin{align}
m^\ast=\frac{g_B^\ast}{g_{B,i}}\left(L^\ast\right)^{-\nu},
\qquad
L^\ast=\frac{Q_\infty(1+r\tau_r)}{\epsilon^\ast},
\label{eq:m_star}
\end{align}
where $L^\ast$ is the stationary population relative to 2023 and $g_B^\ast$ corresponds to the $1.84$ per
cent a year by which output grows at the saturated rate. The required effort combines two gaps: between
the clean progress that decoupling requires and the progress observed, and, with a scale effect, between
today's population and the stationary one. Remediation has the effectiveness $r=2$ of
Section~\ref{sub:illustration}.

\emph{The scale effect.} The scale parameter $\nu$ is not identified by our evidence on fertility, but
outside evidence narrows the range worth considering. Appendix~\ref{app:nu} shows that, across 110 to 150 countries, clean progress
since 2005 does not rise with population: at given income, the elasticity is close to zero, and the
largest value inside a 95 per cent confidence interval is about $0.3$. A national scale effect of $0.5$ or more is therefore hard to reconcile with the data.  A global scale effect is not ruled out, because
clean knowledge spills across borders \citep{aghion2016carbon}, but the evidence that research
productivity falls as research effort grows \citep{jones1995,bloom2020} argues against values near one.
Evidence that innovation responds strongly to market size \citep{acemoglulinn2004} concerns how much
research is done, which in our model is the policy choice $R$, not the productivity of a given research
share. None of this identifies $\nu$, so we treat $0\le\nu\le0.5$ as a sensitivity range: its lower end matches
the cross-country estimates, and its upper end exceeds every national estimate and leaves room for a
global scale effect. We report results at both ends, and also at $\nu=1$.

\begin{table}[h!]\centering
\begin{threeparttable}
\caption{Calibration of the habitat block and clean technology}
\label{tab:recal}
\footnotesize
\begin{tabular}{llll}
\toprule
Parameter & Value & Target or source & Section\\
\midrule
$\alpha_0,\gamma,\eta,\phi_0,\theta,\sigma,\bar g_A,\omega,\nu_h$ & see text & \citet{STRULIK2024104943}; $g_y=1.84\%$ & \ref{sub:illustration}\\
Feedback strength $S$ & $0.209$ & survey estimate $\hat S$ & \ref{sub:xi}\\
Taste share $\varpi$ & $0.54$ & ideal vs realised fertility & \ref{sub:prefs}\\
Taste elasticity $\xi$ & $0.060$ & $\mathcal{K}(\alpha^\ast)\xi=\varpi S$, $\alpha^\ast=0.432$ & \\
Cost elasticity $\chi$ & $0.096$ & $\chi=(1-\varpi)S$ & \\
Pollution elasticity of quality $e_Q$ & $0.235$ & within-country, 1990--2022 & \ref{sub:pollution}\\
Curvature of quality $c_0$ & $4$ & not identified; $1$ and $9$ in Table~\ref{tab:sens} & \\
Large-load elasticity $\mu$ & $0.29$ & $\mu=e_Q(1+c_0)/c_0$ & \\
Country taste level $\alpha_i$ & $\mathcal{A}(\text{TFR}_{2023}/2.1)$ & observed fertility & \ref{sub:xi}\\
Current clean progress $g_{B,i}$ & $(1+g_{\text{int},i})^{-30}-1$ & CO$_2$ per unit of GDP, 2005--2022 & \\
Scale effect $\nu$ & $0$--$0.5$ ($1$) & sensitivity range; cross-country evidence & \ref{app:nu}\\
Quality in tastes $\iota$ & $1$ ($0$) & common index; $e_Q=0.51$ if $\iota=0$ & \ref{sub:split}\\
Remediation effectiveness $r$ & $2$ & Appendix~\ref{app:calib} & \\
\bottomrule
\end{tabular}
\end{threeparttable}
\end{table}

\subsection{Clean-research thresholds by country}
\label{sub:proj}
We compute the required effort from \eqref{eq:m_star} and check it against simulated paths. For the
paths, each country starts in 2023 at its observed fertility. The country level $\alpha_i$
absorbs everything that sets fertility apart from the model's baseline other than habitat, so that the
exercise isolates the habitat feedback: habitat per capita is normalised to one in 2023 and then moves
with population and with the pollution load, which grows with population and income and falls with
clean progress, $g_{B,t}=m\,g_{B,i}L_t^{\nu}$. Generations last thirty years. For simplicity we let income
per capita grow at the saturated rate throughout. By Proposition~\ref{prop:saturation} this is its limit
on a stationary path, and on any decline slow enough for idea growth to saturate, which includes the
decline below the threshold (Proposition~\ref{prop:threshold}); the assumption abstracts from the
transitional slowdown of Remark~\ref{rem:transition}.

\begin{table}[h!]\centering
\begin{threeparttable}
\caption{Clean research needed for a stationary equilibrium, relative to the effort equivalent to current measured clean progress}
\label{tab:thresholds}
\label{tab:gaps}
\footnotesize
\setlength{\tabcolsep}{4.5pt}
\begin{tabular}{lcccccccc}
\toprule
 & TFR & $g_{\text{int}}$ & $\epsilon^\ast$ & $L^\ast$ & \multicolumn{4}{c}{Required effort $m^\ast$}\\
\cmidrule(lr){6-9}
 & 2023 & \% a year & & 2023 $=1$ & $\nu=0$ & $\nu=0.5$ & $\nu=1$\tnote{a} & $\nu=0.5$, $\tau_r=0.1$\\
\midrule
Korea          & 0.72 & $-1.67$ & 37.8 & 0.04 & 1.11 & 5.38 & 26.08 & 4.91\\
China          & 1.00 & $-2.71$ & 16.5 & 0.10 & 0.57 & 1.82 & 5.84 & 1.66\\
Japan          & 1.21 & $-1.73$ & 9.2 & 0.18 & 1.06 & 2.53 & 6.04 & 2.31\\
Italy          & 1.20 & $-2.64$ & 9.4 & 0.17 & 0.59 & 1.43 & 3.47 & 1.31\\
Germany        & 1.44 & $-3.18$ & 5.0 & 0.32 & 0.45 & 0.78 & 1.38 & 0.71\\
Sweden         & 1.43 & $-4.09$ & 5.1 & 0.32 & 0.29 & 0.52 & 0.93 & 0.47\\
United Kingdom & 1.56 & $-4.70$ & 3.7 & 0.44 & 0.22 & 0.34 & 0.51 & 0.31\\
United States  & 1.62 & $-2.79$ & 3.1 & 0.51 & 0.54 & 0.76 & 1.07 & 0.70\\
France         & 1.64 & $-3.08$ & 3.0 & 0.54 & 0.47 & 0.64 & 0.87 & 0.58\\
\bottomrule
\end{tabular}
\begin{tablenotes}\footnotesize
\item $g_{\text{int}}$: trend in CO$_2$ per unit of GDP, 2005--2022 (Global Carbon Budget via Our World in
Data); TFR from UN WPP 2024. $\epsilon^\ast$: habitat per capita, relative to 2023, at which fertility
reaches replacement. $L^\ast=Q_\infty/\epsilon^\ast$: stationary population relative to 2023 when clean
research clears the threshold. $m^\ast$: clean research needed for a stationary population, relative to
the effort that would produce the country's current measured clean progress, from \eqref{eq:m_star};
values above one mean that the current equivalent falls short. Thresholds for a stationary equilibrium;
transitions from 2023 need up to a tenth more (Section~\ref{sub:proj}). Without remediation unless stated.
\item[a] Outside the sensitivity range $0\le\nu\le0.5$ (Section~\ref{sub:recal} and Appendix~\ref{app:nu}).
\end{tablenotes}
\end{threeparttable}
\end{table}

Table~\ref{tab:thresholds} reports the results. Four findings stand out.

\emph{Without a scale effect, current clean progress suffices almost everywhere.} With $\nu=0$, only Korea
and Japan fall short, and narrowly: they need 11 and 6 per cent more clean effort. Every other country
needs at most three-fifths of the effort that produces its current progress. In this case the threshold
depends only on how fast emission intensity falls, and the habitat gap affects the size of the
stationary population but not whether there is one.

\emph{At the upper end of the range, the countries with the lowest fertility fall short.} With $\nu=0.5$,
Korea needs more than five times its current effort equivalent, Japan two and a half times, China nearly twice
and Italy two-fifths more. Germany, the United States, France, Sweden and the United Kingdom clear the
threshold. The ranking follows the stationary population as much as current decoupling. China and Italy
reduce emission intensity almost as fast as the United States, by 2.7 and 2.6 per cent a year against
2.8, but their fertility is so far below replacement that their populations would settle at a tenth and a
sixth of today's, and smaller populations do less clean research. Across the sensitivity range, then,
Korea needs between 1.1 and 5.4 times its current effort equivalent, Japan between 1.1 and 2.5, China between 0.6
and 1.8 and Italy between 0.6 and 1.4. Only at $\nu=1$, outside the sensitivity range, do Germany and the United States also fall
short.

\emph{Remediation helps, but little.} A remediation tax of 10 per cent of labour income raises the
stationary population by a fifth and lowers the required effort by 9 per cent at $\nu=0.5$. That closes
none of the four shortfalls, consistent with Corollary~\ref{cor:threshold}: remediation is an expensive way
to cross the threshold. If quality acts on costs only, $\iota=0$, it is weaker still: the same tax raises
the stationary population by 7 to 8 per cent and lowers the required effort by 3 to 4 per cent. The
required efforts without remediation are 1 to 4 per cent higher than in Table~\ref{tab:thresholds}, for
example $5.59$ for Korea and $2.59$ for Japan at $\nu=0.5$.

\emph{Transitions add to the requirement, and the trap unfolds slowly.} Simulated paths, which include
the undershoot of population on the way to $L^\ast$, require up to a tenth more effort than
\eqref{eq:m_star} in the countries that start closest to replacement. A country below the threshold,
by contrast, does not decline visibly faster at first: at the current effort equivalent, population in 2143 differs
by less than three per cent between $\nu=0.5$ and $\nu=0$. The difference lies in the limit, where a country
below the threshold declines indefinitely, at $0.41$ per cent a year, instead of settling.
Figure~\ref{fig:transitions} shows the paths.

\begin{figure}[h!]\centering
\includegraphics[width=\textwidth]{fig_transitions.pdf}
\caption{Illustrative paths of population and fertility with $\nu=0.5$, at the current effort equivalent
(solid) and at 1.1 times the required effort where the current equivalent falls short (dashed). The timing depends on
assumptions that the evidence does not discipline; Appendix~\ref{app:paths} gives paths for all
countries without a scale effect.}
\label{fig:transitions}
\label{fig:proj}
\end{figure}

\subsection{The stationary population}
\label{sub:statpop}
Where clean research clears the threshold, population settles at $L^\ast$: about half its current level in
France and the United States, a third in Germany and Sweden, a sixth in Japan and Italy, and a few per cent
in Korea. These levels do not depend on the scale effect or on how quickly tastes and costs adjust, although, as
Table~\ref{tab:sens} shows, they depend strongly on the feedback strength and the pollution elasticity.
Because the feedback is weak, $S=0.21$, the habitat gain that fertility needs is large: habitat per capita
must roughly treble in France and the United States, increase ninefold in Japan and Italy, and rise almost
fortyfold in Korea.

\emph{Within the model, a cleaner environment is a substitute for depopulation.} With the common index the
stationary population is $Q_\infty/\epsilon^\ast$: every percentage point of habitat that a cleaner environment
restores is a percentage point of population that need not be lost. If quality acts on costs only, the
conversion is about two-fifths of a point (Proposition~\ref{prop:split}). How large this substitution is depends on the
pollution elasticity of quality and on the maximum attainable gain, which Table~\ref{tab:sens} varies.

\begin{table}[h!]\centering
\begin{threeparttable}
\caption{Long-run population under alternative parameters (per cent of 2023)}
\label{tab:sens}
\footnotesize
\begin{tabular}{lcccccccc}
\toprule
 & $S$ & Korea & China & Japan & Italy & Germany & US & France\\
\midrule
Baseline                                & 0.21 & 4.2 & 9.8 & 17.5 & 17.1 & 32.3 & 51.1 & 53.7\\
Taste only, $\xi=0.146$ (Section~\ref{sub:illustration}) & 0.27 & 11.9 & 21.1 & 31.7 & 31.1 & 49.1 & 68.4 & 70.9\\
Feedback doubled                        & 0.42 & 26.1 & 39.6 & 53.1 & 52.3 & 72.0 & 90.6 & 92.9\\
Pollution elasticity $e_Q=1$            & 0.21 & 19.8 & 45.4 & 81.7 & 79.5 & 150.6 & 238.1 & 250.2\\
Maximum quality gain 39\% ($c_0=1$)     & 0.21 & 3.7 & 8.4 & 15.1 & 14.7 & 27.9 & 44.1 & 46.4\\
Maximum quality gain 83\% ($c_0=9$)     & 0.21 & 4.8 & 11.1 & 19.9 & 19.4 & 36.8 & 58.1 & 61.1\\
Quality in costs only, $\iota=0$         & 0.21 & 3.9 & 9.2 & 16.7 & 16.2 & 31.2 & 49.7 & 52.3\\
\bottomrule
\end{tabular}
\begin{tablenotes}\footnotesize
\item Stationary population $Q_\infty/\epsilon^\ast$ relative to 2023, when clean research clears the
threshold. ``Feedback doubled''
doubles both $\xi$ and $\chi$. The value $e_Q=1$ is illustrative: the cross-country estimates in Table~\ref{tab:pollution} pair a
clear pollution coefficient with a density coefficient close to zero, so they point to values above the
within-country $0.235$ without identifying one; with it the
maximum quality gain is $7.5$-fold. With quality in costs only, $e_Q=0.51$ is recalibrated to the same
reduced form (Section~\ref{sub:split}).
\end{tablenotes}
\end{threeparttable}
\end{table}

The sensitivity analysis sharpens the message. The long-run level of population is governed by two
numbers the evidence pins down only loosely. The feedback strength matters roughly in proportion: doubling
it halves the required increase in log habitat, so that Korea needs about a sixfold rise in habitat per
capita rather than a fortyfold one, and lifts long-run population in Germany, the United States and France to
between seventy and over ninety per cent of today's. The pollution elasticity matters even more. At the
within-country estimate, a cleaner environment offsets only a small part of the gap; at a larger value,
such as $e_Q=1$, which the cross-country estimates do not rule out, ending pollution raises effective habitat enough to stabilise the United States and France at
more than twice their current populations and Japan and Italy at four-fifths. The curvature parameter,
which our data do not identify, moves the results by less than a sixth. Letting quality act on costs only,
with $e_Q$ recalibrated to the same reduced form, lowers long-run population by 3 to 7 per cent, as
Proposition~\ref{prop:split} leads one to expect.

\subsection{What the quantitative results are and are not}
The projections are conditional statements about the habitat channel, not forecasts. They hold fixed
everything that currently sets each country's fertility apart from the model's baseline, they abstract
from migration and from age structure, and they apply to national aggregates a feedback estimated from
variation across settlements. The cross-section also tells us nothing about how quickly tastes and costs
adjust when habitat changes over time; Section~\ref{sub:caveats} notes that the habitat channel accounts
for little of the observed decline in fertility. Reading current clean progress from the trend in emission
intensity treats all of that trend as the equivalent of clean research, and the scale effect is not
identified; we report a sensitivity range. Read with these qualifications, the exercise delivers a quantitative answer to the question
the paper poses. The habitat mechanism makes a stationary population the long-run outcome only where
clean research clears a threshold, and where fertility is lowest that threshold is high, while the
population at which it settles is small.

% ============================================================
\section{Policy}
\label{sec:policy}
% ============================================================

The model has two policy instruments, and Corollary~\ref{cor:threshold} assigns them distinct roles.
Clean research decides whether population settles at all; remediation decides, given that it settles, at
what level. Without a scale effect the division is complete. Both roles presume that clean progress is
proportional to research effort; if clean knowledge is semi-endogenous, neither instrument secures a
stationary population (Proposition~\ref{prop:semiendog}). With a scale effect, remediation also supports
decoupling, because the larger population it sustains does more clean research, but the effect is
weak: with $r=2$ and $\nu=0.5$, remediation would be the cheaper route across the threshold only if the
threshold share exceeded the whole labour force. Neither instrument changes the long-run growth rate,
and both lower the level of income per capita. Clean research withdraws labour from production, and
both instruments, being financed from labour income, reduce education spending, which the
human-capital technology amplifies: a one per cent cut in education inputs lowers human capital by
$\theta/(1-\theta)\approx5$ per cent in the long run.

Whether these costs are worth bearing depends on how the planner values population. Remediation is the
simpler case, because with clean research above the threshold it affects only the level of the
stationary population and income. Appendix~\ref{sub:welfare} shows that a total-utilitarian
planner, who values the additional lives that remediation sustains, chooses a positive remediation tax at
our baseline, with a value of life calibrated to the value of a statistical life, whereas an
average-utilitarian planner, at the baseline effectiveness of remediation, does so only if environmental
quality is valued directly and the tax spares education spending. With the common
habitat index, the gains in habitat per capita that remediation buys are absorbed in the long run by a
larger population. If quality acts on costs only, less is absorbed, and the two criteria move apart: the
total-utilitarian tax falls and the average-utilitarian tax rises. When clean
research falls just short of the threshold, a modest remediation tax can carry the economy across it, but
the decline it prevents lies more than a thousand years ahead and, even at a discount rate of a quarter of a
per cent a year, adds little to the case for the tax.

For clean research the population criterion decides even more (Appendix~\ref{sub:welfare_R}). An
average-utilitarian planner who places no direct value on the environment chooses no clean research at
all, and accepts an indefinite decline. At a discount rate of one per cent a year, a total-utilitarian planner chooses two and a half to four
times the threshold share at the model baseline, and about twice it at a quarter of a per cent, because clean research both secures a stationary population and speeds the
recovery of habitat. The exception is the economy furthest from replacement. At the upper end of the
scale effect, Korea must devote nearly 8 per cent of labour to clean research to cross the threshold, and
because the tax that pays for it also lowers human capital, every criterion we consider, at either discount rate,
prefers to stay below it. Two further margins are left aside. Households take habitat per capita as given, so their fertility crowds
later generations, and a planner with a fertility instrument would combine it with environmental policy,
as \citet{gerlagh2023} show for a carbon tax and a tax on births; remediation, which raises
fertility along the transition, has a natalist effect, although through habitat rather than, as in
\citet{DELACROIX2012271}, through a tax that makes procreation relatively cheaper. And
under the bounded ideas law there is no long-run scale effect of population on income; if there were one,
as \citet{edenkuruc2026} argue, the larger population that both instruments sustain would raise income per
head, which would strengthen the case for both.

% ============================================================
\section{Conclusion}
\label{sec:conclusion}
% ============================================================

Idea-based growth models with endogenous fertility leave open what could stop a decline in population.
\citet{jones22} shows that a shrinking population can stall idea production; \citet{STRULIK2024104943}
shows that growth can continue when ideas are produced by human capital, and his result applies at a
constant population too. Neither supplies a force that brings fertility back to replacement. This paper
proposes one. When the taste for children and the time needed to raise them depend on habitat per capita,
a declining population raises the space and environmental quality available to each person, fertility
rises, and replacement becomes a unique and locally stable outcome rather than a knife-edge.

Whether the economy gets there depends on technology. Output keeps growing, and the habitat shrinks
unless clean knowledge grows faster than output. When clean knowledge comes from research, this
condition becomes a threshold for clean research, which, since clean knowledge is assumed to earn no
private return, the market does not supply. When clean innovation rises with the number of people, the
threshold also depends on demography: the economies with the lowest fertility settle at the smallest
populations and need the most clean research, and below the threshold population decline slows clean
innovation until the environment produces an Empty Planet of its own. In simulations, an economy can
fall into this trap during the undershoot of population even when the stationary state would be
sustainable once reached. Remediation of environmental quality raises the size of the stationary
population; it supports decoupling only through the scale effect, and weakly. Neither clean research
nor remediation changes the long-run growth rate, which is set by the ceiling of Strulik's bounded law.
Whether either is worth its cost depends on how population is valued: an average-utilitarian planner
chooses neither unless the environment is valued directly, whereas a total-utilitarian planner chooses
both, except where fertility is lowest and the scale effect strong.

The evidence is consistent with the central premise, although it is associational. Fertility is
positively associated with proxies for habitat per capita across settlements in the World Values Survey,
in the European Values Study, in two US surveys and across 161 countries, and under the maintained
conversion from settlement size to habitat the implied feedback lies well inside the stable region. A
pure time-cost account is rejected within the maintained labour-supply specification, and stated ideals
rise with habitat at about half the rate of realised fertility, consistent with a role for preferences.
At a given density and income, fertility is also higher where the pollution load on the habitat is
lower, although reverse causality from fertility to emissions cannot be excluded.

Calibrated to these estimates, the model gives a sobering quantitative answer. Over a sensitivity range
for the scale effect informed by cross-country evidence, Korea needs between 1.1 and 5.4 times the
clean effort that produces its current measured clean progress, and Japan between 1.1 and 2.5; China and
Italy fall short only at the upper end of the range, and the other economies we consider clear the
threshold throughout. Where clean research clears the threshold, population settles only after large
declines, and the lower the fertility, the larger the decline: at the within-country pollution
elasticity, to about half its present level in France and the United States, a sixth in Japan and
Italy, and a few per cent in Korea. A larger pollution elasticity would raise these levels substantially.

These conclusions depend on how clean knowledge is produced. With clean progress proportional to research
effort, a stationary population is a policy outcome, and the thresholds above say what it takes. If clean
knowledge is semi-endogenous, with diminishing returns to its own stock, a permanent research share
raises its level but not its growth rate, and at the returns estimated for ideas in general no share
secures a stationary population (Proposition~\ref{prop:semiendog}). Whether population settles is then a
question of technology that policy cannot answer. We regard the two cases as equally important
alternatives; the evidence on clean innovation does not yet discriminate between them.

Three extensions seem most valuable. The first is the source of clean research: directed technical change
driven by an emissions price, as in \citet{beulmann2025}, would make the threshold a condition on that
price and connect our results to the fertility-to-emissions channel they study. The second is the scale
effect and the returns to clean research, which our results show to be decisive and for which we report
only a sensitivity range; evidence on how clean innovation scales with research effort at the global
level would narrow it and discriminate between the fully and semi-endogenous laws. The third is
empirical: direct measures of local environmental quality matched to fertility data, and the time series
of countries in which population decline has begun, would test the pollution channel and the stabilising
dynamics that the cross-section can only suggest.

\appendix
\setcounter{table}{0}\renewcommand{\thetable}{\Alph{section}.\arabic{table}}
% ============================================================
\section{Separating tastes from costs through household budgets}
\label{app:budget}

%% -------------------------------------------------------------
\subsection{What the model restricts, and what it does not}
\label{sub:restrictions}
%% -------------------------------------------------------------

A negative association between settlement size and fertility is among the most robustly documented
regularities in demography, and it is implied by essentially every model in the literature we build
on. In \citet{DAO2013442} density operates on the \emph{cost} of children, and in \citet{dao25} population
density raises the \emph{time cost} of children: in our notation, both raise $\phi$. In \citet{dao25}
resource depletion works differently, raising fertility by widening the gender gap in education. In \citet{DELACROIX201713} the same reduced-form
correlation is documented without committing to a channel. Our model adds a channel of a different kind:
habitat also enters \emph{preferences}, through $\alpha$, and not only the budget set. Establishing that fertility
falls with settlement size is therefore consistent with our model but provides no evidence for it
over these alternatives.

The closed forms of Proposition~\ref{prop:hab_solution} nonetheless separate the channels, though
not in the way a first glance suggests. Labour supply is the natural place to look, since a taste
shift moves it while a pure time-cost shift does not. That comparison is not robust, however, to
the cost the density literature emphasises most. Housing enters as a per-child \emph{goods} cost,
and a goods cost moves labour supply in the same direction as a taste shift. Within the model, what
separates tastes from every cost is a different object: under logarithmic preferences the
household's expenditure shares are pinned down by $\alpha$ alone, and cost parameters of any kind
redistribute resources between fertility and per-child inputs without disturbing them. The result
is exact, but it inherits the unit income elasticities of those preferences, and
Section~\ref{sub:ltest} shows that this is where it meets the data least well.

\begin{proposition}[Expenditure shares identify the taste for children]
\label{prop:identification}
Extend the household problem of Section~\ref{sec:model} to allow a per-child goods cost
$\kappa_t\ge 0$ and a price $p_t>0$ for education goods, so that \eqref{eq:budget_hab} becomes
\begin{align}
w_t h_t\left(1-\phi n_t-s_t n_t\right)=c_t+\left(p_t q_t+\kappa_t\right)n_t .
\label{eq:budget_general}
\end{align}
Let $\alpha_t>\Gamma$ and write $W_t\equiv w_t h_t$. Then the consumption share, total education
time and total education spending share depend on $\alpha_t$ alone,
\begin{align}
\frac{c_t}{W_t}=\frac{1}{1+\alpha_t},
\qquad
s_t n_t=\frac{\gamma\eta}{1+\alpha_t},
\qquad
\frac{p_t q_t n_t}{W_t}=\frac{\gamma\theta}{1+\alpha_t},
\label{eq:invariants}
\end{align}
and are invariant to $\phi$, $\kappa_t$ and $p_t$; fertility is
\begin{align}
n_t=\frac{(\alpha_t-\Gamma)\,W_t}{(1+\alpha_t)\left(\phi W_t+\kappa_t\right)} ;
\label{eq:n_general}
\end{align}
and labour supply depends on the \emph{composition} of the cost of children,
\begin{align}
\ell_t=\frac{1+\alpha_t-\gamma\eta-\psi_t\left(\alpha_t-\Gamma\right)}{1+\alpha_t},
\qquad
\psi_t\equiv\frac{\phi W_t}{\phi W_t+\kappa_t}\in(0,1],
\label{eq:ell_general}
\end{align}
which reduces to \eqref{eq:sol_l_hab} when $\kappa_t=0$.
\end{proposition}

\begin{proof}
The first-order conditions for $q_t$ and $s_t$ give $p_tq_tn_t=\gamma\theta c_t$ and
$W_ts_tn_t=\gamma\eta c_t$. That for $n_t$ is
$\alpha_t/n_t=\left[W_t(\phi+s_t)+p_tq_t+\kappa_t\right]/c_t$, which on substituting the previous
two becomes $(\alpha_t-\Gamma)/n_t=(\phi W_t+\kappa_t)/c_t$, i.e.
$(\phi W_t+\kappa_t)n_t=(\alpha_t-\Gamma)c_t$. Substituting all three into
\eqref{eq:budget_general} gives
$c_t(1+\Gamma)=W_t-(\phi W_t+\kappa_t)n_t=W_t-(\alpha_t-\Gamma)c_t$, hence
$c_t=W_t/(1+\alpha_t)$. The remaining expressions follow, the last from
$\ell_t=1-\phi n_t-s_tn_t$ with $\phi n_t=\psi_t(\alpha_t-\Gamma)/(1+\alpha_t)$.
\end{proof}

Two implications organise what follows. First, a shock to the price of education goods is not a
rival explanation at all: $p_t$ enters neither \eqref{eq:n_general} nor \eqref{eq:ell_general}, so
it cannot generate a fertility gradient, and only $q_t$ adjusts. The rivals are $\phi$ and
$\kappa_t$. Second, because each quantity in \eqref{eq:invariants} is proportional to
$1/(1+\alpha_t)$ while fertility is not, the ratio of their habitat elasticities eliminates
$\partial\ln\alpha/\partial\ln\epsilon$ and leaves a prediction in preference parameters only:
\begin{align}
\Lambda'\;\equiv\;
\frac{\partial\ln\left(pqn/W\right)/\partial\ln\epsilon}{\partial\ln n/\partial\ln\epsilon}
=-\frac{\alpha-\Gamma}{1+\Gamma},
\label{eq:Lambda}
\end{align}
with the same value if the numerator is replaced by $sn$ or by $c/W$. At our calibration ---
$\Gamma=0.234$ and $\alpha$ between $\alpha_0=0.39$ and $\alpha^\ast=0.4316$ ---
\begin{align}
\Lambda'\in[-0.160,\,-0.126],
\label{eq:Lambda_number}
\end{align}
whereas $\Lambda'=0$ under \emph{every} cost channel, since the numerator of \eqref{eq:Lambda} is
then identically zero. The corresponding ratio built on labour supply,
\begin{align}
\Lambda_\ell=\frac{\alpha-\Gamma}{1+\Gamma}
\left[\frac{(1-\psi)(1+\alpha)}{1+\alpha-\gamma\eta-\psi(\alpha-\Gamma)}-1\right],
\label{eq:Lambda_ell}
\end{align}
coincides with $\Lambda'$ when $\psi=1$ but is contaminated otherwise. Table~\ref{tab:fingerprints}
collects the comparison; Section~\ref{sub:ltest} takes both ratios to the data, beginning with
$\Lambda_\ell$, which requires the fewest auxiliary assumptions.

\begin{table}[h!]\centering
\begin{threeparttable}
\caption{How the channels differ}
\label{tab:fingerprints}
\begin{tabular}{lcccc}
\toprule
Response to a roomier & Taste & Time cost & Goods cost & Education price\\
habitat & ($\alpha\uparrow$) & ($\phi\downarrow$) & ($\kappa\downarrow$) & ($p\downarrow$)\\
\midrule
Fertility $n$                            & $+$ & $+$ & $+$ & $0$\\
Education per child $s$, $q$             & $-$ & $-$ & $-$ & $q$ rises\\
Education totals $sn$, $pqn/W$           & $-$ & $0$ & $0$ & $0$\\
Consumption share $c/W$                  & $-$ & $0$ & $0$ & $0$\\
Labour supply $\ell$                     & $-$ & $0$ & $-$ & $0$\\
\midrule
$\Lambda'$, eq.~\eqref{eq:Lambda}        & $[-0.160,-0.126]$ & $0$ & $0$ & n.a.\\
$\Lambda_\ell$, eq.~\eqref{eq:Lambda_ell}& eq.~\eqref{eq:Lambda_ell} & $0$ & $(-0.128,0)$ & n.a.\\
\bottomrule
\end{tabular}
\begin{tablenotes}\footnotesize
\item Signs from Proposition~\ref{prop:identification}. A fall in $p$ leaves fertility, labour
supply and all three invariants unchanged and raises $q$ one for one, so it generates no habitat
gradient and the ratios are not defined. Under the taste channel $\Lambda_\ell$ equals $\Lambda'$
only when $\psi=1$; under a goods-cost channel it sweeps $(-0.128,0)$ as $\psi$ falls. The ratio
$s/q^{\ast}=\eta/\theta$ is invariant to every parameter here and serves as a specification check.
\end{tablenotes}
\end{threeparttable}
\end{table}

\paragraph{Two channels at once.}
When the time cost of Section~\ref{sec:model} responds to habitat as well, the ratios take on a
second role. Under the two-channel model the habitat elasticities of fertility and of the education
spending share are
\begin{align}
\pi_n=\mathcal{K}(\alpha)\,\xi+\chi,
\qquad
\pi_Q=-\frac{\alpha}{1+\alpha}\,\xi,
\label{eq:two_elasticities}
\end{align}
so that the ratio \eqref{eq:Lambda} becomes
\begin{align}
\Lambda'=\frac{\pi_Q}{\pi_n}
=-\frac{\alpha-\Gamma}{1+\Gamma}\cdot\frac{\mathcal{K}\xi}{\mathcal{K}\xi+\chi} .
\label{eq:Lambda_share}
\end{align}
The first factor is the pure-taste value in \eqref{eq:Lambda}; the second is the \emph{share of the
fertility gradient attributable to tastes}. The band \eqref{eq:Lambda_number} is thus the
prediction of the taste channel acting alone, $\Lambda'=0$ that of costs acting alone, and an
estimate between the two measures the mix. The two reduced-form coefficients identify both
structural elasticities,
\begin{align}
\xi=-\frac{1+\alpha}{\alpha}\,\pi_Q,
\qquad
\chi=\pi_n+\frac{1+\Gamma}{\alpha-\Gamma}\,\pi_Q,
\label{eq:decomposition}
\end{align}
and with them the feedback strength $S$ and the speed of adjustment of
Proposition~\ref{prop:two_channels}.

%% -------------------------------------------------------------
\subsection{Separating the channels: labour supply and expenditure shares}
\label{sub:ltest}
%% -------------------------------------------------------------

Proposition~\ref{prop:identification} turns the question into one of sign and ratio. Two ratios
are available. The first pairs fertility with labour supply; it separates a pure time cost from
everything else and needs only hours worked. The second pairs fertility with an expenditure share;
within the model it separates tastes from every cost, but it requires budget data and, as will be
seen, more of the preference specification than the data grant. We set out the second in full,
since the first is the same construction with hours in the numerator. Fertility and the education
spending share are regressed on the same habitat measure over a common sample and with an
identical right-hand side:
\begin{align}
\ln n_{it} &= \pi_n \ln \epsilon_{it} + \mathbf{x}_{it}'\boldsymbol{\delta}_n+\iota_t+u_{it},
\label{eq:pair_n}\\
\ln\!\left(p_{it}q_{it}n_{it}/W_{it}\right) &= \pi_{Q}\ln \epsilon_{it}
  + \mathbf{x}_{it}'\boldsymbol{\delta}_{Q}+\iota_t+\upsilon_{it}.
\label{eq:pair_q}
\end{align}
The model predicts
\begin{align}
\pi_n>0,
\qquad
\pi_{Q}<0,
\qquad
\widehat\Lambda'=\hat\pi_{Q}/\hat\pi_n\in[-0.160,\,-0.126],
\label{eq:test}
\end{align}
with the confidence interval for $\widehat\Lambda'$ obtained by the delta method or by a cluster
bootstrap. Every cost channel --- a time cost $\phi$ that rises with crowding \citep{DAO2013442,dao25} or with environmental
degradation, a housing or other per-child goods cost $\kappa$, and the price of education goods
--- predicts $\pi_{Q}=0$ and hence $\Lambda'=0$, because by \eqref{eq:invariants} none of them
touches the expenditure shares. The hypotheses are therefore nested and separately testable, and
failing to reject $\pi_{Q}=0$ is informative against our mechanism rather than merely
uninformative. Either of the other two invariants in \eqref{eq:invariants} --- total education time
$sn$, or the consumption share $c/W$ --- may replace the numerator, with the same predicted value.
If the time cost of children also varies with habitat, as in Section~\ref{sec:model}, an
estimate between zero and the band in \eqref{eq:test} is not a rejection of either mechanism: by
\eqref{eq:Lambda_share} it measures the share of the fertility gradient that operates through
tastes.

\paragraph{Is the test powered?}
The restriction in \eqref{eq:test} is a ratio of two small numbers, and its precision is governed by
the dispersion of habitat per capita in the sample, which within a single country is small.
Appendix~\ref{app:power} simulates the test under the sampling design of the General Social Survey.
With the calibrated fertility gradient a sample of that size would detect the taste channel four times
in five; at the gradient actually observed in the GSS, which is half as steep, power falls below one in
three, and a decisive test needs of the order of $150{,}000$ to $600{,}000$ observations.

\paragraph{Labour supply.}
With hours worked in the numerator of \eqref{eq:test} the ratio is $\widehat\Lambda_\ell$. It
discriminates less finely but is far more widely measured and does not depend on how spending
responds to income. By \eqref{eq:ell_general} it equals $\Lambda'$ when the cost of
children is purely a time cost, and departs from it as the goods share rises: under the taste
channel it takes the value \eqref{eq:Lambda_ell}, while under a pure goods-cost channel it sweeps
$(-0.128,\,0)$. A labour-supply test therefore still rejects a pure time-cost channel, such as
that of \citet{dao25}, for which $\Lambda_\ell=0$ exactly, but cannot on its
own separate a taste shift from a housing-cost shift. We take it to two US surveys: the General Social Survey, which measures individual hours alongside
fertility and size of place over five decades, and the Consumer Expenditure Survey, whose
estimation sample is about 1.7 times larger and which records the hours of both spouses.

\paragraph{Results from the General Social Survey.}
We begin with the GSS cumulative file, 1975--2022, for respondents aged
25--54. Habitat is minus the log of size of place (\texttt{SIZE}, in thousands); fertility is
children ever born (\texttt{CHILDS}); labour supply is weekly hours (\texttt{HRS1}) with
non-workers entered at zero, as the model's $\ell$ requires. Both equations are estimated by
Poisson pseudo-maximum likelihood on a common sample, with year effects and controls for age, sex,
marital status, race, region and religious attendance, post-stratification weights, and a
bootstrap over the $6{,}184$ sampling clusters. Table~\ref{tab:gss} in Appendix~\ref{app:power}
reports the headline estimate and the subsamples discussed below.



Two findings emerge, and they pull in different directions. The fertility gradient replicates in
an independent survey of a single country: $\hat\pi_n$ is positive and precisely estimated in every
row, and at $\varrho=0.85$ the headline implies a feedback strength of $\hat S=0.162$
$[0.127,\,0.197]$, somewhat below the $0.209$ obtained from the World Values Survey but of the same
order. Labour supply, by contrast, shows no detectable response to habitat. The headline ratio is
$-0.046$ with an interval that contains both zero and the whole of the taste-only band, so the test
discriminates between neither hypothesis --- as Table~\ref{tab:power} leads one to expect at this
sample size. The sign is as the taste channel predicts for women and in the first half of the
sample, where the ratios of $-0.31$ and $-0.14$ are close to or beyond the band, and it is reversed
for men and after 2000. On the intensive margin the evidence is sharper and unfavourable: among
those in work the ratio is $0.011$ with an interval of $[-0.070,\,0.083]$ that excludes the band,
so any response of labour supply to habitat operates through participation and not through hours.
Adding education and income raises the ratio instead of attenuating it; with a headline
indistinguishable from zero that movement cannot be read as the bad-control bias described below.
We also note that the ideal number of children (\texttt{CHLDIDEL}) shows no gradient in size of
place in the GSS, with an elasticity of $0.0018$ (s.e.\ $0.0015$) at ages 25--54 and $0.0003$ at
ages 25--64 with childhood habitat controlled (Table~\ref{tab:prefs}), in contrast to the World
Values Survey.

We conclude that the labour-supply test is not informative at the scale of the GSS. It does not
reject the time-cost channel, and it does not reject the taste channel.

\paragraph{Results from the Consumer Expenditure Survey.}
Our second source is the Interview Survey of the US Consumer Expenditure Survey, data releases 2000--2025 (interviews conducted 2000--2026), which
supports both ratios: $662{,}987$ quarterly interviews, which we average within consumer
unit to obtain $236{,}529$ households. The estimation sample is the $64{,}853$ units with a
reference person aged 25--44, one or two adults, and positive earnings. Three features of the
implementation matter. First, Proposition~\ref{prop:identification} defines the shares relative to
\emph{full} income $W$, not to spending: relative to consumption, education outlays are the
constant $\gamma\theta$ whatever the habitat. We therefore scale each numerator by the household's
hourly wage, family wage and salary income divided by the annual hours of the reference person and
spouse. Second, the public files report location only as the population class of the sampling
area, five bands plus a rural indicator; habitat is minus the log of the class midpoint, and
inference is by bootstrap over 139 state-by-class cells because sampling-unit identifiers are
withheld for most households. Third, fertility is measured by children under 18 present in the
household. All equations include year effects and controls for age, sex, marital status, race and
region, and are weighted.

\begin{table}[h!]\centering
\begin{threeparttable}
\caption{Expenditure-share and labour-supply ratios in the Consumer Expenditure Survey}
\label{tab:cex}
\footnotesize\setlength{\tabcolsep}{4pt}
\begin{tabular}{lccc}
\toprule
Outcome & Habitat elasticity & Ratio to $\hat\pi_n$ & 95\% interval\\
\midrule
\multicolumn{4}{l}{\emph{A. Headline: earners, $N=64{,}853$}}\\
Children under 18, $\hat\pi_n$                    & $\phantom{-}0.0370$ $(0.0053)$ & & \\
Education spending / wage                         & $-0.0129$ $(0.0166)$ & $-0.35$ & $[-1.29,\,0.43]$\\
Any education spending                            & $\phantom{-}0.0036$ $(0.0052)$ & $\phantom{-}0.10$ & $[-0.21,\,0.34]$\\
Consumption / wage                                & $\phantom{-}0.0088$ $(0.0024)$ & $\phantom{-}0.24$ & $[\phantom{-}0.13,\,0.34]$\\
Household annual hours                            & $-0.0123$ $(0.0021)$ & $-0.33$ & $[-0.45,\,-0.24]$\\
Hourly wage                                       & $-0.0586$ $(0.0060)$ & & \\
\midrule
\multicolumn{4}{l}{\emph{B. Conditioning on the log wage}}\\
Children under 18, $\hat\pi_n$                    & $\phantom{-}0.0326$ $(0.0054)$ & & \\
Education spending / wage                         & $-0.0259$ $(0.0160)$ & $-0.80$ & $[-1.94,\,0.07]$\\
Consumption / wage                                & $-0.0229$ $(0.0032)$ & $-0.70$ & $[-1.08,\,-0.53]$\\
Household annual hours                            & $-0.0129$ $(0.0023)$ & $-0.39$ & $[-0.55,\,-0.28]$\\
\midrule
\multicolumn{4}{l}{\emph{C. All one- and two-adult households, $N=74{,}902$}}\\
Children under 18, $\hat\pi_n$                    & $\phantom{-}0.0353$ $(0.0049)$ & & \\
Household annual hours, non-workers at zero       & $-0.0103$ $(0.0018)$ & $-0.29$ & $[-0.44,\,-0.20]$\\
\bottomrule
\end{tabular}
\begin{tablenotes}\footnotesize
\item Elasticities with respect to minus the log population of the sampling area. Education
spending is winsorised at the 99th percentile and estimated by Poisson pseudo-maximum likelihood,
as are the counts and the hours in panel C; consumption, hours in panels A and B, and the wage are
in logs. Consumption is total expenditure less education. Bootstrap standard errors in
parentheses and percentile intervals, 200 replications. Tastes alone predict a ratio in
$[-0.160,-0.126]$ for the shares; any cost channel predicts $0$.
\end{tablenotes}
\end{threeparttable}
\end{table}

The fertility gradient appears for a third time, in a third survey: $\hat\pi_n=0.037$, or
$\hat S=0.25$ $[0.18,\,0.32]$ at $\varrho=0.85$, close to the World Values Survey estimate.

Labour supply responds, and robustly. Household hours fall with habitat in every specification,
with a ratio between $-0.29$ and $-0.39$ and intervals that exclude zero. Within the maintained
labour-supply specification this rejects the pure time-cost channel, for which $\Lambda_\ell=0$, in a
sample large enough to have power against it,
and it has the sign that both the taste channel and a goods-cost channel predict. It is also
larger than either can deliver, which points to the violation discussed under the qualifications
below: wages are about six per cent lower for each log point by which the population of the sampling
area is smaller, and hours respond to wages directly.
The GSS, where no response was detectable, sits uneasily beside this: its interval reaches
$-0.20$, which only touches the upper end of the intervals obtained here.

\paragraph{Expenditure shares: a restriction the data reject.}
The share test fares differently. The education share, the numerator on which
Section~\ref{sub:restrictions} places most weight, is uninformative: fewer than a third of these
households report any education outlay while in the survey, the amounts are highly skewed and
include adults' own tuition, and the ratio of $-0.35$ carries an interval, $[-1.29,\,0.43]$, that
contains zero, the taste band and a great deal else.

The consumption share is precisely estimated but behaves as neither hypothesis requires.
Unconditionally it \emph{rises} with habitat, giving a ratio of $0.24$ that rejects both the taste
band and zero. Conditional on the wage it falls, with a ratio of $-0.70$ that has the sign of the
taste channel and several times its magnitude. The reversal identifies the problem. Consumption
relative to the wage declines steeply in the wage itself, whereas the logarithmic preferences
behind Proposition~\ref{prop:identification} impose a unit income elasticity and hence shares that
are independent of $W$. That invariance fails in the data, and the estimated ratio takes whatever
sign the treatment of the wage gives it.

\paragraph{Relaxing homotheticity.}
The failure is one of functional form, and the most economical repair keeps the rest of the model.
Let utility from consumption be $\log(c_t-\bar c)$ with a subsistence level $\bar c>0$. The
first-order conditions are those of Proposition~\ref{prop:identification} with $c_t$ replaced by
$c_t-\bar c$, and the budget constraint then gives
\begin{align}
c_t-\bar c=\frac{W_t-\bar c}{1+\alpha_t},
\qquad
p_tq_tn_t=\frac{\gamma\theta\left(W_t-\bar c\right)}{1+\alpha_t},
\qquad
\left(\phi W_t+\kappa_t\right)n_t=\frac{\left(\alpha_t-\Gamma\right)\left(W_t-\bar c\right)}{1+\alpha_t}.
\label{eq:stone_geary}
\end{align}
The consumption share now falls with income, as in the data, and since $\bar c/W_t\to0$ along a
growth path the long-run results of Sections~\ref{sec:model}--\ref{sec:theory} are
unaffected. What depends on $\alpha_t$ alone is no longer the share but the \emph{slope} of
consumption, or of education spending, in full income: $1/(1+\alpha_t)$ and
$\gamma\theta/(1+\alpha_t)$. Cost parameters leave both slopes unchanged, and the ratio of the
habitat elasticity of either slope to that of fertility has the same taste-only value
\eqref{eq:Lambda}.



Table~\ref{tab:sg} in Appendix~\ref{app:power} reports the result. In the headline sample the ratio is
$-2.19$ $[-3.30,\,-1.03]$ for the consumption slope and $-3.18$ $[-5.94,\,-0.74]$ for the education
slope. The slopes fall with habitat, as the taste channel
requires; the cost-only prediction of zero is rejected for the consumption slope in the full sample
and in both halves of it, and for the education slope in the full sample and the first half; and
the consumption and education slopes both move in the predicted direction. Their point
estimates differ, $-2.19$ against $-3.18$ in the headline sample, but given the wide intervals they are
not statistically distinguishable, as the proportionality that \eqref{eq:stone_geary} imposes
requires. The magnitudes, however, are an
order of magnitude beyond the taste band. Something other than tastes is therefore also moving the
slopes, and the leading candidates are local prices, which differ systematically with density, and
error in measured wages, whose attenuating effect on an estimated slope need not be the same in
large and small labour markets. We read Table~\ref{tab:sg} as qualitative support for a taste
component and as no basis for quantifying it.

\paragraph{Where this leaves the channels.}
The positive association between fertility and habitat is robust across four surveys and a
cross-country panel.
A pure time-cost account of it is rejected on the labour-supply margin, within the maintained
labour-supply specification. Between tastes and a goods
cost the budget-based evidence is suggestive and no more: spending slopes move in the direction
that tastes predict, but by too much to be attributed to them alone.
Section~\ref{sub:prefs} turns to evidence that does not pass through the budget constraint.
For the purposes of this paper the distinction matters less than it might appear. By
Proposition~\ref{prop:two_channels} and Corollary~\ref{cor:goods} population stabilises whichever of the
taste, time-cost and goods-cost channels is active, and the speed at which it does so depends on the feedback strength $S$ and not
on its composition.

\paragraph{A by-product: the composition of the cost of children.}
Because $\Lambda'$ is invariant to $\psi$ and $\Lambda_\ell$ is not, estimating both identifies the
time share of the cost of children. The cleanest route is through levels rather than elasticities:
$s_tn_t$ is observable as total time devoted to children's education and is equal to
$\gamma\eta/(1+\alpha_t)$ by \eqref{eq:invariants}, so $\phi n_t=1-\ell_t-s_tn_t$ and
\begin{align}
\psi_t=\frac{\left(1-\ell_t-s_tn_t\right)\left(1+\alpha_t\right)}{\alpha_t-\Gamma},
\label{eq:psi_identified}
\end{align}
with $\alpha_t$ recovered from any of the invariants. The share of the cost of children borne as
time rather than goods is a quantity the density literature argues over and has not pinned down;
our framework delivers it in principle as a by-product of the identification exercise. In practice
it inherits the dependence on the share invariance documented above, and we do not report an
estimate.

\paragraph{Qualifications.}
Two apply to both tests. The model has homogeneous adults, whereas in the data the labour-supply
margin is overwhelmingly female; we therefore estimate the labour-supply ratio separately by sex and
treat the household-level prediction as the one the theory speaks to. And the ratios are joint
tests of the mechanism and of the exclusion restriction that settlement size affects the numerator
only through the fertility taste. The most plausible violation is that wages differ with density, and the sensitivity of the
consumption ratio to conditioning on the wage shows that it matters. A third qualification is
specific to the model: $s$, $q$ and labour income are chosen jointly with $n$ and $\ell$ given
$\alpha$, so conditioning on education or on labour income conditions on an outcome of the
mechanism and biases both ratios toward zero. The headline specifications omit both, and we report
the attenuation as controls are added rather than suppressing it.


%% -------------------------------------------------------------
\subsection{Power of the ratio test and further estimates}
\label{app:power}

\paragraph{Power of the ratio test.}
The restriction in \eqref{eq:test} is a ratio of two small numbers, and what governs its precision
is the dispersion of habitat per capita in the sample, not of settlement size. In the General
Social Survey used in Appendix~\ref{sub:ltest} the standard deviation of log size of place is $2.18$, which at
$\varrho=0.85$ corresponds to a standard deviation of $\ln\epsilon$ of only $0.33$, three-quarters of
it between primary sampling units. We simulate the two-channel model of
Section~\ref{sec:model} from its closed forms under exactly that design: six respondents per
cluster, children drawn as Poisson counts, a cluster shock with standard deviation $0.20$ common to
both equations, and a numerator whose idiosyncratic log standard deviation of $0.9$ is calibrated
to the sampling error of weekly hours in the GSS. The simulated ratio is therefore the labour-supply
ratio $\Lambda_\ell$ with $\psi=1$, which then coincides with $\Lambda'$. The null world has the time cost alone,
calibrated to the same fertility gradient as the alternative, so the two are observationally
identical in the fertility equation.

\begin{table}[h!]\centering
\begin{threeparttable}
\caption{Power of the ratio test under a realistic sampling design}
\label{tab:power}
\footnotesize\setlength{\tabcolsep}{3.5pt}
\begin{tabular}{lrccccc}
\toprule
& & \multicolumn{2}{c}{Cost-only null} & \multicolumn{3}{c}{Taste channel active}\\
\cmidrule(lr){3-4}\cmidrule(lr){5-7}
Fertility gradient & $N$ & s.d. & size & true & mean (s.d.) & power\\
\midrule
Calibrated; tastes only, $S=0.32$           & $37{,}800$  & $0.048$ & $0.04$ & $-0.126$ & $-0.127$ $(0.053)$ & $0.82$\\
                                           & $150{,}000$ & $0.025$ & $0.04$ & $-0.126$ & $-0.129$ $(0.024)$ & $1.00$\\
Calibrated; tastes + time cost, $\chi=0.1$, $S=0.42$ & $37{,}800$ & $0.048$ & $0.04$ & $-0.097$ & $-0.100$ $(0.040)$ & $0.71$\\
                                           & $150{,}000$ & $0.025$ & $0.04$ & $-0.097$ & $-0.097$ $(0.019)$ & $1.00$\\
Observed in GSS; tastes only, $S=0.16$  & $37{,}800$ & $0.109$ & $0.07$ & $-0.126$ & $-0.128$ $(0.099)$ & $0.30$\\
                                           & $150{,}000$ & $0.052$ & $0.04$ & $-0.126$ & $-0.123$ $(0.053)$ & $0.77$\\
                                           & $600{,}000$ & $0.024$ & $0.06$ & $-0.126$ & $-0.124$ $(0.024)$ & $0.99$\\
\bottomrule
\end{tabular}
\begin{tablenotes}\footnotesize
\item 200 replications per row. ``Size'' and ``power'' refer to a one-sided test of $\Lambda_\ell=0$ at
the five per cent level against the simulated null distribution. $N=37{,}800$ approximately matches the GSS
estimation sample. $S$ is the habitat elasticity of fertility in \eqref{eq:S_def}; the calibrated
value $0.32$ is $\mathcal{K}\xi$ evaluated at $\alpha_0$, where the simulation is centred, whereas
Section~\ref{sub:illustration} evaluates it at $\alpha^\ast$, giving $0.275$. Adding a time-cost
elasticity $\chi=0.1$ gives $0.42$.
\end{tablenotes}
\end{threeparttable}
\end{table}

The ratio is estimated without bias in every configuration, and with the fertility gradient of the
calibration a sample of GSS size would detect the taste channel four times in five. The gradient
actually observed in the GSS is half as steep, however, and precision deteriorates in proportion:
the standard error of the ratio doubles to about $0.10$ and power falls to below one in three. A
decisive test at the observed gradient needs of the order of $150{,}000$ to $600{,}000$ observations, or a sample
with more dispersion in habitat per capita than a single country provides.

\begin{table}[h!]\centering
\begin{threeparttable}
\caption{The labour-supply ratio $\widehat\Lambda_\ell$ in the General Social Survey}
\label{tab:gss}
\footnotesize\setlength{\tabcolsep}{3.5pt}
\begin{tabular}{lrcccc}
\toprule
Sample and specification & $N$ & $\hat\pi_n$ & $\hat\pi_\ell$ & $\widehat\Lambda_\ell$ & 95\% interval\\
\midrule
Headline                              & $37{,}678$ & $0.0243$ $(0.0027)$ & $-0.0011$ $(0.0022)$ & $-0.046$ & $[-0.204,\,0.142]$\\
\quad + education                     & $37{,}614$ & $0.0190$ $(0.0024)$ & $\phantom{-}0.0038$ $(0.0022)$ & $\phantom{-}0.199$ & $[-0.038,\,0.451]$\\
\quad + education and income          & $34{,}923$ & $0.0196$ $(0.0025)$ & $\phantom{-}0.0053$ $(0.0023)$ & $\phantom{-}0.269$ & $[\phantom{-}0.060,\,0.531]$\\
Women                                 & $20{,}806$ & $0.0204$ $(0.0039)$ & $-0.0064$ $(0.0038)$ & $-0.311$ & $[-0.652,\,0.070]$\\
Men                                   & $16{,}872$ & $0.0300$ $(0.0041)$ & $\phantom{-}0.0028$ $(0.0026)$ & $\phantom{-}0.092$ & $[-0.064,\,0.261]$\\
Married, household hours              & $20{,}916$ & $0.0189$ $(0.0026)$ & $\phantom{-}0.0023$ $(0.0018)$ & $\phantom{-}0.120$ & $[-0.080,\,0.323]$\\
Workers only                          & $28{,}289$ & $0.0285$ $(0.0030)$ & $\phantom{-}0.0003$ $(0.0011)$ & $\phantom{-}0.011$ & $[-0.070,\,0.083]$\\
1975--1998                            & $19{,}389$ & $0.0262$ $(0.0033)$ & $-0.0036$ $(0.0028)$ & $-0.136$ & $[-0.331,\,0.068]$\\
2000--2022                            & $18{,}289$ & $0.0205$ $(0.0043)$ & $\phantom{-}0.0016$ $(0.0032)$ & $\phantom{-}0.078$ & $[-0.211,\,0.418]$\\
\bottomrule
\end{tabular}
\begin{tablenotes}\footnotesize
\item Coefficients are elasticities with respect to minus log size of place; multiply by
$1/(1-\varrho)$ to express them per unit of $\ln\epsilon$. Bootstrap standard errors in parentheses
and percentile intervals, 200 replications over sampling clusters. The taste channel alone predicts
$\Lambda_\ell\in[-0.160,-0.126]$ when $\psi=1$; a pure time cost predicts $0$.
\end{tablenotes}
\end{threeparttable}
\end{table}

\paragraph{Slopes of spending in full income.}
Table~\ref{tab:sg} reports the Stone--Geary slope ratios discussed in Appendix~\ref{sub:ltest}.

\begin{table}[h!]\centering
\begin{threeparttable}
\caption{Slopes of spending in full income, Consumer Expenditure Survey}
\label{tab:sg}
\footnotesize\setlength{\tabcolsep}{4pt}
\begin{tabular}{lcccc}
\toprule
& \multicolumn{2}{c}{Consumption slope} & \multicolumn{2}{c}{Education slope}\\
\cmidrule(lr){2-3}\cmidrule(lr){4-5}
Sample & ratio & 95\% interval & ratio & 95\% interval\\
\midrule
Headline, $N=64{,}853$ & $-2.19$ & $[-3.30,\,-1.03]$ & $-3.18$ & $[-5.94,\,-0.74]$\\
Married                & $-0.89$ & $[-2.02,\,\phantom{-}0.07]$ & $-1.14$ & $[-4.24,\,\phantom{-}2.05]$\\
2000--2012             & $-3.21$ & $[-5.78,\,-1.29]$ & $-5.39$ & $[-9.38,\,-1.05]$\\
2013 onwards           & $-1.48$ & $[-2.28,\,-0.40]$ & $-2.10$ & $[-5.26,\,\phantom{-}1.93]$\\
Wage between 8 and 80  & $-1.59$ & $[-2.31,\,-0.64]$ & $-1.85$ & $[-4.06,\,\phantom{-}1.16]$\\
\bottomrule
\end{tabular}
\begin{tablenotes}\footnotesize
\item Annual consumption and education spending per adult are regressed on full income (the hourly
wage times $2{,}000$), habitat and their interaction, with the controls of Table~\ref{tab:cex}.
Each entry is the habitat elasticity of the income slope divided by $\hat\pi_n$. Tastes alone
predict $[-0.160,-0.126]$; any cost channel predicts $0$. Percentile intervals from 200 bootstrap
replications over state-by-class cells.
\end{tablenotes}
\end{threeparttable}
\end{table}

\clearpage
\section{A cross-country consistency check}
\label{app:crosscountry}\label{sub:crosscountry}
\setcounter{table}{0}

This appendix reports a check that requires no micro data. Proposition~\ref{prop:inversion} maps an
observed fertility rate into the taste parameter that rationalises it, and
$\alpha(\epsilon)=\alpha_0\epsilon^{\xi}$ then maps that taste into the habitat per capita the model
requires. Running both steps on observed national fertility therefore produces, for each country, a
\emph{predicted population density} --- a quantity the model was not designed to match, and one
that can be compared with the observed figure.

The exercise has exactly one free normalisation, which we state explicitly. The model's own baseline
is $\epsilon=1$, at which $\alpha=\alpha_0=0.39$ and hence $n=0.813$: with replacement at a total
fertility rate of $2.1$, that corresponds to a TFR of $1.71$. We set the density that
corresponds to $\epsilon=1$ to $D_0=120.5$ inhabitants per square kilometre, the observed density of
France, whose TFR of $1.64$ is the closest in our sample to the model's baseline. This fixes the scale
only: France itself is not assigned $\epsilon=1$, and its own fertility implies $\hat\epsilon=0.883$.
Nothing else is chosen: $\xi=0.146$, $\phi$ and
$\Gamma$ are all carried over unchanged from Section~\ref{sub:illustration}, and every other entry in
Table~\ref{tab:crosscountry} follows mechanically.

\begin{table}[h!]\centering
\begin{threeparttable}
\caption{Model-implied and observed population density, 2023}
\label{tab:crosscountry}
\begin{tabular}{lcccrrcc}
\toprule
& & & & & & \multicolumn{2}{c}{observed / implied}\\
\cmidrule(lr){7-8}
& TFR & $\hat\alpha$ & $\hat\epsilon$ & $D_0/\hat\epsilon$ & observed $D$ & $\xi=0.146$ & $\xi=0.123$\\
\midrule
Korea       & $0.72$ & $0.295$ & $0.149$ & $810$ & $523$  & $0.65$ & $0.45$\\
Taiwan      & $0.87$ & $0.309$ & $0.202$ & $596$ & $658$  & $1.11$ & $0.82$\\
Japan       & $1.21$ & $0.340$ & $0.394$ & $306$ & $330$  & $1.08$ & $0.91$\\
Italy       & $1.20$ & $0.339$ & $0.386$ & $312$ & $201$  & $0.64$ & $0.54$\\
France      & $1.64$ & $0.383$ & $0.883$ & $137$ & $120$  & $0.88$ & $0.86$\\
Singapore   & $0.94$ & $0.316$ & $0.234$ & $514$ & $8{,}476$ & $16.48$ & $12.56$\\
\bottomrule
\end{tabular}
\begin{tablenotes}\footnotesize
\item Total fertility and population density (inhabitants per square kilometre) are 2023 estimates
from UN World Population Prospects 2024. $\hat\alpha=\mathcal{A}(\text{TFR}/2.1)$ from
\eqref{eq:inversion}; $\hat\epsilon=(\hat\alpha/\alpha_0)^{1/\xi}$, shown for the calibrated
$\xi=0.146$. France's observed density supplies the scale normalisation $D_0=120.5$. The last column repeats the comparison at
the survey estimate \eqref{eq:xi_estimate}.
\end{tablenotes}
\end{threeparttable}
\end{table}

The model reproduces the broad pattern but not every pairwise ranking: it places Korea above Taiwan and Italy above Japan, where observed densities rank them the other way. The magnitudes are close. Taiwan and Japan are matched to
within about ten per cent and France to within twelve; Korea and Italy are the two cases where the
observed density falls short of what the model requires, meaning that their fertility is lower than
the habitat channel alone accounts for. At the smaller survey estimate $\xi=0.123$ the model needs
more crowding to rationalise a given fertility shortfall, and every ratio falls: Taiwan, Japan and
France remain within twenty per cent, while Korea and Italy move to roughly half. For a one-parameter preference shifter whose elasticity was
fixed in Section~\ref{sub:illustration} for an unrelated reason, agreement of this order across a
fivefold range of densities among the countries other than Singapore is more than the exercise was entitled to expect.

Two cautions keep us from claiming more. Singapore is off by a factor of sixteen, as any city-state
must be: crude density is not measuring habitat where there is no hinterland. And the result is
sensitive to the land concept. Crude density happens to fit well, but replacing it with habitable
land --- arguably the better proxy for $E_t$ --- would raise the observed density of the mountainous
cases of Korea and Japan several-fold, and turn Japan's close agreement and Korea's shortfall into
substantial overshooting. That sensitivity, and the fact that the survey estimate of $\xi$ in the last
column is itself drawn from the relationship between fertility and settlement size, make Table~\ref{tab:crosscountry} an illustrative internal consistency
check rather than an independent validation.

\paragraph{The residual: levels, not changes.}
The check above is cross-sectional. Its time-series counterpart does not pass. Korea's
population rose from 26 to 52 million between 1960 and 2023, halving habitat per capita; at
$\xi=0.146$ this moves the model's fertility from a TFR of $1.71$ to $1.33$, whereas Korean
fertility fell from $6.0$ to
$0.72$. The habitat channel can therefore account for why Korea sits far below other high-income
countries at a point in time, but for only a modest part of its decline over time. This is not a
refutation --- the model never claimed to explain the demographic transition, most of which operates
through forces held fixed here --- but it does mean the channel is partial. It also leaves the
central dynamic prediction untested: stabilisation at $\epsilon^{\ast}$ requires population decline
to raise $\epsilon$, and in no country has that yet begun on a scale that would show up in fertility.
The cross-section can show that the mechanism is present; only time can show that it stabilises.


%\clearpage
\clearpage
\section{Welfare analysis of remediation and clean research}
\label{sub:welfare}
\setcounter{table}{0}
This appendix asks when remediation raises welfare and, in Appendix~\ref{sub:welfare_R}, how much
clean research a planner would choose. We use the full model of
Section~\ref{sec:model}, with the parameters of Table~\ref{tab:recal}: clean knowledge follows
\eqref{eq:B_law} with $\nu=0.5$, the upper end of the sensitivity range of Section~\ref{sub:recal}, and clean research, financed by a tax on labour
income, is held at $1.25$ times its threshold $R^\ast(0)$ of Corollary~\ref{cor:threshold}, so that
remediation changes the level of the stationary population but not whether there is one. With the scaling
of Section~\ref{sub:illustration} and $\nu=0$, this share gives clean knowledge the growth that decoupling
at $\zeta=7$ would give. Appendix~\ref{app:calib} varies $\nu$ and also places clean research just below
the threshold, where remediation decides whether there is a stationary population. In this appendix
$\tau$ denotes the remediation tax $\tau_r$ alone, and the common habitat index, $\iota=1$, applies unless
stated. Two choices have to be made. The first concerns population: with endogenous fertility, a
policy that changes the number of people must be ranked by a criterion that takes a view on the value of
additional lives \citep{golosov2007}. We use the average-utilitarian (Millian) criterion as the main
case, the discounted sum of period utility of a representative member of each generation,
$W^{M}=\sum_t\delta^tu_t$, and report the total-utilitarian (Benthamite) criterion
$W^{B}=\sum_t\delta^tL_t(u_t+\bar u)$ as the alternative, where $\delta$ is the generational discount
factor and $\bar u$ is the utility level at which a life is just worth living. The second concerns the taste for children, which habitat itself moves. Our
main case evaluates each generation with the taste it actually has, $\alpha(\epsilon_t)$; the
alternative evaluates every allocation with the same reference taste $\alpha^\ast$, so that a policy is
not credited with changing preferences. In both cases we allow environmental quality to enter utility
directly, as an amenity, through a term $\Psi\ln Q_t$ that households take as given and that therefore
leaves their choices unchanged; $\Psi=0$ is the model of the previous sections.

\begin{proposition}[Remediation in the long run]
\label{prop:welfare}
Let habitat enter tastes and costs through the common index, $\iota=1$, let a constant clean-research
share $R$ clear the threshold of Proposition~\ref{prop:threshold}, so that the economy settles at
$L^\ast(0)$ without remediation, let $Q_\infty(\tau)$ be long-run quality, and compare
stationary-population paths with and without a remediation tax $\tau$ for a given path of ideas. Then:
\begin{enumerate}
\item[(i)] habitat per capita, the taste for children and fertility are the same,
$\epsilon=\epsilon^\ast$, $\alpha=\alpha^\ast$ and $n=1$, and the stationary population is proportional
to $Q_\infty(\tau)$;
\item[(ii)] if the tax is levied on labour income, period utility differs by
\begin{align}
\Delta u(\tau)=\frac{1+\gamma\theta}{1-\theta}\ln\frac{1-R-\tau}{1-R}+\Psi\ln\frac{Q_\infty(\tau)}{Q_\infty(0)},
\label{eq:du_labour}
\end{align}
and if the revenue is instead raised by a tax on consumption at the rate $\tau_c$, with
$\tau_cc_t=\tau W_t\ell_t$, which spares education spending, by
\begin{align*}
\Delta u(\tau)=\ln\left[1-\frac{(1+\gamma\theta)\tau}{1-R}\right]+\Psi\ln\frac{Q_\infty(\tau)}{Q_\infty(0)};
\end{align*}
\item[(iii)] under the Millian criterion a small tax raises long-run welfare if and only if
\begin{align*}
\Psi\,\frac{Q_\infty'(0)}{Q_\infty(0)}&>\frac{1+\gamma\theta}{(1-\theta)(1-R)}
&&\text{with labour-income financing, or}\\
\Psi\,\frac{Q_\infty'(0)}{Q_\infty(0)}&>\frac{1+\gamma\theta}{1-R}
&&\text{with education exempt.}
\end{align*}
Under the Benthamite criterion, applied to the stationary path from a given date, it does so if and only
if $\left[Q_\infty'(0)/Q_\infty(0)\right]\left(\bar U+\Psi\right)$ exceeds the same cost, where
$\bar U\equiv(1-\delta)\sum_t\delta^t(u_t+\bar u)$ is the discounted average of period utility plus $\bar u$
along the stationary path without the tax. $\bar U$ is finite although $u_t$ grows without bound,
because $u_t$ grows linearly in $t$ at the saturated rate.
\end{enumerate}
\end{proposition}

\begin{proof}
Part (i) is Proposition~\ref{prop:decoupling}(i). Since $L^\ast(\tau)\ge L^\ast(0)$, the threshold
condition holds under any tax if it holds without one, habitat per capita converges to $\epsilon^\ast$
whatever the tax, and $L^\ast=\bar E\,Q_\infty(\tau)/\epsilon^\ast$. Since $n=1$, the term $\alpha\ln b$ equals
$\alpha^\ast\ln2$ under every tax and either treatment of tastes. With labour-income financing, consumption and education goods scale
with $1-R-\tau$, and by \eqref{eq:h_tax_dyn} human capital on the stationary path scales with
$(1-R-\tau)^{\theta/(1-\theta)}$; collecting the terms in $\ln c_t+\gamma\ln h_{t+1}$ gives the coefficient
$1+\theta/(1-\theta)+\gamma\theta/(1-\theta)=(1+\gamma\theta)/(1-\theta)$. With a consumption tax the budget
constraint is $(1-R)W_t(1-\phi n_t-s_tn_t)=(1+\tau_c)c_t+q_tn_t$, with $W_t=w_th_t$. In terms of
$\tilde c_t\equiv(1+\tau_c)c_t$ the household problem is the untaxed one, because utility changes only by
the constant $-\ln(1+\tau_c)$; so $n_t$, $s_t$ and $q_t$ are unchanged, human capital is unaffected and
$\tilde c_t=(1-R)W_t/(1+\alpha^\ast)$. The revenue condition $\tau_c c_t=\tau W_t\ell(\alpha^\ast)$, with
$(1+\alpha^\ast)\ell(\alpha^\ast)=1+\gamma\theta$, gives $\tau_c/(1+\tau_c)=\tau(1+\gamma\theta)/(1-R)$ and hence
$c_t=W_t\left[1-R-\tau(1+\gamma\theta)\right]/(1+\alpha^\ast)$. Part (iii) follows by differentiating
$\Delta u$ and $\ln\big[L^\ast\sum_t\delta^t(u_t+\bar u+\Delta u)\big]$ at $\tau=0$.
\end{proof}

The proposition makes precise what Corollary~\ref{cor:threshold} implies for welfare. In the
long run remediation converts environmental quality into population: habitat per capita, and with it
the taste for children, returns to $\epsilon^\ast$. An average-utilitarian planner therefore gains from
remediation only through the direct value of a cleaner environment, and must weigh it against a cost
that the human-capital technology amplifies: a labour-income tax also cuts education spending, and a
one per cent cut in education inputs lowers human capital by $\theta/(1-\theta)\approx5$ per cent in the
long run. A total-utilitarian planner also values the larger population, and may favour remediation even
when the environment has no direct value.

Table~\ref{tab:welfare} quantifies these results along the full transition of the re-calibrated model,
from the model's baseline (a total fertility rate of $1.71$) and from Korea's starting point, with
$\nu=0.5$, clean research at $1.25$ times its threshold in both cases, an annual discount rate of one per
cent, and the remediation technology of Section~\ref{sub:illustration}, $Q_\infty(\tau)\propto1+2\tau$.
Because Korea's stationary population is small, its threshold is high: clean research takes $2.5$ per cent
of effective labour at the baseline and $9.7$ per cent in Korea. For the Benthamite criterion we follow \citet{adhami2024}, who calibrate the intercept of log flow
utility to a value of a statistical life of \$7.4 million, which puts flow utility at $4.87$ in units
where consumption is one. We set $\bar u$ so that the full 2023 utility index of our model plus
$\bar u$, $\ln c+\alpha\ln b+\gamma\ln h+\Psi\ln Q+\bar u$, equals $4.87$, with 2023 consumption
normalised to one and $\ln Q=0$ in 2023. The terms in children and human capital are thus absorbed into
the calibrated intercept rather than added to it, and $\bar u$ itself differs from $4.87$ by those terms. Appendix~\ref{app:calib} discusses this and the
other parameters that the data do not identify.

\begin{table}[h!]\centering
\begin{threeparttable}
\caption{Optimal remediation tax}
\label{tab:welfare}
\footnotesize
\setlength{\tabcolsep}{5pt}
\begin{tabular}{llcccccc}
\toprule
 & & \multicolumn{2}{c}{Millian, actual tastes} & \multicolumn{2}{c}{Millian, fixed tastes} & \multicolumn{2}{c}{Benthamite, actual tastes}\\
\cmidrule(lr){3-4}\cmidrule(lr){5-6}\cmidrule(lr){7-8}
Financing & Amenity $\Psi$ & $\tau^\ast$ & gain & $\tau^\ast$ & gain & $\tau^\ast$ & gain\\
\midrule
\multicolumn{8}{l}{\emph{A. Model baseline, TFR $1.71$}}\\
Labour income   & 0   & 0    & 0    & 0    & 0    & 0.09 & 1.1\\
                & 0.5 & 0    & 0    & 0    & 0    & 0.17 & 3.7\\
                & 1   & 0    & 0    & 0    & 0    & 0.25 & 7.8\\
Education exempt& 0   & 0    & 0    & 0    & 0    & 0.40 & 11.5\\
                & 0.5 & 0    & 0    & 0    & 0    & 0.47 & 18.6\\
                & 1   & 0.07 & 1.0  & 0.07 & 1.0  & 0.52 & 26.4\\
\midrule
\multicolumn{8}{l}{\emph{B. Korea, TFR $0.72$}}\\
Labour income   & 0   & 0    & 0    & 0    & 0    & 0.03 & 0.1\\
                & 0.5 & 0    & 0    & 0    & 0    & 0.16 & 1.9\\
                & 1   & 0    & 0    & 0    & 0    & 0.25 & 5.7\\
Education exempt& 0   & 0    & 0    & 0    & 0    & 0.16 & 1.3\\
                & 0.5 & 0    & 0    & 0    & 0    & 0.29 & 5.2\\
                & 1   & 0    & 0    & 0    & 0    & 0.37 & 10.6\\
\bottomrule
\end{tabular}
\begin{tablenotes}\footnotesize
\item Optimal permanent tax on a grid of one percentage point between $0$ and $0.6$. Millian gains are consumption equivalents in
per cent; Benthamite gains are per cent changes in $W^{B}$. Re-calibrated model of Table~\ref{tab:recal},
$\nu=0.5$, clean research at $1.25$ times its threshold $R^\ast(0)$, financed by a labour-income tax, annual
discount rate one per cent, $Q_\infty(\tau)\propto1+2\tau$, flow utility of $4.87$ in 2023. ``Education exempt'': a
consumption tax raising the same revenue (Proposition~\ref{prop:welfare}(ii)). The Benthamite criterion
evaluates utility with actual tastes.
Appendix~\ref{app:calib} reports the optimal tax for other values of the remediation technology, of
$\bar u$ and of the curvature $c_0$, and for remediation that abates the pollution load instead of
restoring quality.
\end{tablenotes}
\end{threeparttable}
\end{table}

Three results stand out. First, under the average-utilitarian criterion a remediation tax is not optimal
unless households value environmental quality directly and the tax spares education spending: with
labour-income financing even an amenity weight equal to that of consumption, $\Psi=1$, does not justify
it, because the direct gain, $\Psi\,Q_\infty'/Q_\infty=2$ at $\tau=0$, falls short of the long-run cost of
$7.3$. In Korea the same planner chooses no tax even with $\Psi=1$ and education exempt, although the long-run
condition of Proposition~\ref{prop:welfare} holds there, $2$ against $(1+\gamma\theta)/(1-R)=1.35$. The condition compares
stationary paths. From Korea's starting point the gain in quality arrives only as pressure fades, over
many generations, while the cost is paid from the first, and at a discount rate of one per cent the cost
dominates. Second, the treatment of tastes hardly matters. The habitat channel moves fertility only along the
transition, and in the long run fertility is at replacement under every tax. Third, the population
criterion matters a great deal. A total utilitarian values the population that remediation sustains and
chooses a positive tax in every case. With labour-income financing and no amenity, the tax is nine per
cent at the model baseline and three per cent in Korea, where clean research already absorbs a tenth of
labour income. When education is exempt and the environment has amenity value, it exceeds half at the
model baseline and is almost forty per cent in Korea. Remediation is thus, in this model, primarily a
population policy, and whether it is worth having depends on how one values the people it brings into
existence.

\subsection{Parameters the data do not identify}
\label{app:calib}

Besides the scale effect $\nu$, which Section~\ref{sec:quant} varies throughout and which we vary at the
end of this appendix, three parameters are not identified by the evidence of Section~\ref{sec:empirics}: the effectiveness of remediation, the largest gain in environmental quality
that ending pollution could deliver, and the utility level at which a life is just worth living. This
appendix explains the values chosen, anchors them where outside evidence exists, and reports how the
results change when each is varied, and when the discount rate is lowered. Table~\ref{tab:calib_sens} collects the welfare results; all other
settings are those of Table~\ref{tab:welfare}, panel~A.

\paragraph{Effectiveness of remediation.}
We model remediation as restoring environmental quality, $Q_\infty(\tau)\propto1+r\tau$, so that
$r=\mathrm{d}\ln Q_\infty/\mathrm{d}\tau$ at $\tau=0$: spending one per cent of output permanently raises
long-run quality by $r$ per cent. We know of no estimate of this return for restoration. The closest
evidence concerns abatement, which lowers the pollution load rather than restoring quality. In the
cost-effective mitigation pathways assessed by \citet{riahi2022} that limit warming to $2^\circ$C
($>67$ per cent), net greenhouse-gas emissions fall by 52--76 per cent between 2019 and 2050, and global
GDP in 2050 is 1.3--2.7 per cent below its reference level, an estimate that excludes the benefits of
avoided climate damages. These figures refer to all greenhouse gases. Carbon dioxide, our measure of the
load, falls at least as fast in the same pathways, by half in the 2040s and to net zero in the 2070s,
so applying the greenhouse-gas reductions to the carbon-dioxide load, if anything, understates the
effectiveness of abatement. At the midpoints, a 64 per cent cut in the load for 2 per cent of output
raises quality in our specification by about 23 per cent on impact, a return of roughly $r\approx11$. The
comparison is only a heuristic, because it sets a cost in one year against a permanent expenditure, and
we do not read it as evidence that $r=2$ is conservative; it suggests only that $r=2$ is not implausibly
large. Abatement differs in one essential
respect, however: under decoupling the load vanishes in the long run whatever the policy, so abatement
speeds the transition but does not change long-run quality or the stationary population.

Proposition~\ref{prop:welfare} makes the role of $r$ transparent: a small tax raises long-run
average-utilitarian welfare if and only if $\Psi r$ exceeds $7.33$ with labour-income financing, or $1.25$
with education exempt, at the baseline clean-research share of $2.5$ per cent. Panel~A of Table~\ref{tab:calib_sens} confirms this along the transition. With
$\Psi=1$, the average-utilitarian planner chooses a sizeable tax only once $r$ approaches $10$ with
labour-income financing ($0.01$ at $r=5$ and $0.09$ at $r=10$), or $2$ with education exempt; without an
amenity value it never does. The
total-utilitarian optimum rises steadily with $r$. Panel~B replaces restoration by abatement, with the
abatement technology calibrated to the IPCC midpoints, $Q=\big[(1+c_0)/(1+c_0e^{-b\tau}x)\big]^{\mu}$ with
$b=-\ln(0.36)/0.02=51$. Optimal taxes are then small, at most eight per cent, whether clean research
lies above or just below the threshold, because the benefit is concentrated on the transition, but they
are positive under the average-utilitarian criterion once the environment has amenity value.

\paragraph{The largest gain in quality.}
The curvature $c_0$ in $Q=\big[(1+c_0)/(1+c_0x)\big]^{\mu}$ measures the current load relative to the
scale at which pollution begins to bite, and fixes the largest attainable gain in quality,
$(1+c_0)^{\mu}$. If that scale is taken to be the country's 1950 emission level, $c_0$ lies between $0.6$
and $2$ for the United Kingdom, Germany, Sweden, France and the United States, between $8$ and $10$ for Japan and
Italy, and above $100$ for China and Korea. Our baseline $c_0=4$ and the range $1$--$9$ in
Table~\ref{tab:sens} cover the advanced economies. The parameter matters for the projections, by about
fourteen per cent either way, but hardly for welfare: in Proposition~\ref{prop:welfare} only
$Q_\infty'/Q_\infty=r$ enters, and the factor $(1+c_0)^{\mu}$ cancels. Panel~C of
Table~\ref{tab:calib_sens} confirms that the optimal tax changes by at most two percentage points.

\paragraph{The value of a life.}
The total-utilitarian criterion requires a utility level at which a life is just worth living. We take it
from the literature that calibrates the intercept of log flow utility to the value of a statistical life
\citep{jonesklenow2016,adhami2024}. \citet{adhami2024} set $\bar u=4.87$, in units where US consumption per person in 2006 is
one: a value of a statistical life of \$7.4 million in 2006 prices \citep[the figure used by the US
Environmental Protection Agency; see][]{adhami2024}, spread over a remaining life expectancy of about 40
years, values a year of life at \$185{,}000, or $4.87$ times consumption per person of \$38{,}000.
\citet{jonesklenow2016} obtain $\bar u=5.00$ from a value of \$6 million in 2007 prices. With log flow
utility, both imply that flow utility is zero at about 0.7--0.8 per cent of US consumption per person. We use $4.87$ as the 2023 value of flow utility, and vary it to $2.4$ and $7.3$, the values implied by a value of life half
as large and half as large again. Panel~D of Table~\ref{tab:calib_sens}
shows that this parameter matters for the size of the total-utilitarian optimum but, with one exception,
not for its sign: only with labour-income financing, no amenity value and the low value of life does the
optimal tax fall to zero. The average-utilitarian results do not depend on $\bar u$.

\paragraph{The scale effect.}
Panel~E varies $\nu$. With clean research at $1.25$ times its threshold, the scale effect hardly matters:
the optimal tax changes by at most one percentage point between $\nu=0$ and $\nu=1$, because
Proposition~\ref{prop:welfare} holds for any $\nu$ once the threshold is cleared. The panel also places
clean research five per cent below its threshold. Without remediation, population then declines
indefinitely, and by Corollary~\ref{cor:threshold}(a) a remediation tax lowers the threshold. In the
simulations a tax of six per cent at $\nu=0.5$, or four per cent at $\nu=1$, carries the economy across it
and restores a stationary population. Yet the optimal tax rises by at most two percentage points. The
decline that remediation prevents sets in only after more than a thousand years: without remediation,
population at $\nu=0.5$ is still $0.57$ of its 2023 level after forty generations, and that far future
carries little weight.

\paragraph{The discount rate.}
Panel~F lowers the annual discount rate from one per cent to a half and a quarter of a per cent, with a
horizon of 200 generations. A lower discount rate raises the optimal tax, because the gains from
remediation, a larger population and, with an amenity value, a cleaner environment, last forever, while
the cost of the transition comes early. With education exempt, the total-utilitarian optimum rises from
$0.40$ to $0.60$, the upper end of the grid, and the average-utilitarian optimum with $\Psi=1$ from $0.07$ to $0.20$; with
labour-income financing the average-utilitarian optimum stays at zero. The rows with clean research below
the threshold differ from those above it by at most two percentage points even at a quarter of a per cent.
Within the range of discount rates used in long-run analysis, then, preventing an indefinite decline that
begins in the distant future adds little to the case for remediation.

\paragraph{Quality in costs only.}
Panel~G lets quality act on the costs of children only, $\iota=0$, with $e_Q$ recalibrated as in
Section~\ref{sub:recal}. Proposition~\ref{prop:welfare}(i) no longer applies: by
Proposition~\ref{prop:split}, remediation now raises the stationary population by only about two-fifths of
the gain in quality, and the rest raises habitat per capita and changes long-run tastes. The two criteria
move in opposite directions. The total-utilitarian optimum falls, to zero with labour-income financing,
because remediation buys less population. The average-utilitarian optimum with an amenity value rises,
from $0.07$ to $0.13$ at the model baseline, because less of the cleaner environment is absorbed by a
larger population. With education exempt, both criteria then choose a modest tax at the model baseline.

\begin{table}[h!]\centering
\begin{threeparttable}
\caption{Optimal remediation tax under alternative values of the unidentified parameters}
\label{tab:calib_sens}
\scriptsize
\setlength{\tabcolsep}{3pt}
\begin{tabular}{lcccc}
\toprule
 & \multicolumn{2}{c}{Labour-income financing} & \multicolumn{2}{c}{Education exempt}\\
\cmidrule(lr){2-3}\cmidrule(lr){4-5}
 & Millian, $\Psi=1$ & Benthamite, $\Psi=0$ & Millian, $\Psi=1$ & Benthamite, $\Psi=0$\\
\midrule
\multicolumn{5}{l}{\emph{A. Effectiveness of restoration, $r$}}\\
$r=0.5$ & 0 & 0 & 0 & 0\\
$r=1$ & 0 & 0 & 0 & 0.23\\
$r=2$ (baseline) & 0 & 0.09 & 0.07 & 0.40\\
$r=5$ & 0.01 & 0.23 & 0.23 & 0.52\\
$r=10$ & 0.09 & 0.27 & 0.29 & 0.56\\
\midrule
\multicolumn{5}{l}{\emph{B. Abatement instead of restoration, $b=51$ (IPCC midpoints)}}\\
Clean research $1.25\,R^\ast$ & 0.03 & 0.05 & 0.06 & 0.08\\
Clean research $0.95\,R^\ast$ (below threshold) & 0.03 & 0.06 & 0.06 & 0.08\\
\midrule
\multicolumn{5}{l}{\emph{C. Curvature of quality, $c_0$}}\\
$c_0=1$ & 0 & 0.08 & & \\
$c_0=4$ (baseline) & 0 & 0.09 & & \\
$c_0=9$ & 0 & 0.11 & & \\
\midrule
\multicolumn{5}{l}{\emph{D. Value of a life: 2023 flow utility (Benthamite, $\Psi=0$ and $\Psi=1$)}}\\
$2.4$ (VSL halved) & \multicolumn{2}{c}{0 \;/\; 0.14} & \multicolumn{2}{c}{0.25 \;/\; 0.44}\\
$4.87$ (baseline) & \multicolumn{2}{c}{0.09 \;/\; 0.25} & \multicolumn{2}{c}{0.40 \;/\; 0.52}\\
$7.3$ (VSL $\times1.5$) & \multicolumn{2}{c}{0.21 \;/\; 0.35} & \multicolumn{2}{c}{0.49 \;/\; 0.57}\\
\midrule
\multicolumn{5}{l}{\emph{E. Scale effect $\nu$ and the level of clean research}}\\
$\nu=0$, $1.25\,R^\ast$ & 0 & 0.09 & 0.06 & 0.39\\
$\nu=0.5$, $1.25\,R^\ast$ (baseline) & 0 & 0.09 & 0.07 & 0.40\\
$\nu=1$, $1.25\,R^\ast$ & 0 & 0.09 & 0.07 & 0.40\\
$\nu=0.5$, $0.95\,R^\ast$ (below threshold) & 0 & 0.11 & 0.07 & 0.40\\
$\nu=1$, $0.95\,R^\ast$ (below threshold) & 0 & 0.10 & 0.08 & 0.41\\
\midrule
\multicolumn{5}{l}{\emph{F. Annual discount rate; clean research $1.25\,R^\ast$ / $0.95\,R^\ast$}}\\
1 per cent (baseline) & 0 / 0 & 0.09 / 0.11 & 0.07 / 0.07 & 0.40 / 0.40\\
0.5 per cent & 0 / 0 & 0.11 / 0.12 & 0.12 / 0.12 & 0.58 / 0.59\\
0.25 per cent & 0 / 0 & 0.12 / 0.14 & 0.20 / 0.21 & 0.60 / 0.60\\
\midrule
\multicolumn{5}{l}{\emph{G. Quality in costs only, $\iota=0$, $e_Q=0.51$}}\\
Model baseline & 0 & 0 & 0.13 & 0.13\\
Korea & 0 & 0 & 0.09 & 0\\
\bottomrule
\end{tabular}
\begin{tablenotes}\footnotesize
\item Optimal permanent tax on a grid of one percentage point between $0$ and $0.6$, model baseline,
$\nu=0.5$ and clean research at $1.25$ times its threshold $R^\ast(0)$ unless stated, annual discount rate
one per cent. Entries of $0.60$ lie at the upper end of the grid. Panel~F uses a horizon of 200 generations. Benthamite entries use actual tastes. Without an amenity value
($\Psi=0$) the Millian optimum is zero in every case. Panel~C reports labour-income financing only. 
\end{tablenotes}
\end{threeparttable}
\end{table}

\subsection{The optimal clean-research share}
\label{sub:welfare_R}
We now hold remediation at zero and let the planner choose a permanent clean-research share $R$, financed
by a labour-income tax as in Section~\ref{sec:model}. Unlike remediation, clean research decides whether
there is a stationary population, so the criteria now compare a stationary population with an indefinite
decline. Proposition~\ref{prop:welfare} still describes the cost: each point of $R$ lowers long-run
period utility as a labour-income tax does, by $(1+\gamma\theta)/[(1-\theta)(1-R)]$ per unit, because it
also lowers education spending. Above the threshold, more clean research does not change the stationary
population or long-run quality; it only shortens the time that pressure takes to fade. Table~\ref{tab:welfare_R}
reports the optimal share relative to the threshold, on a grid of one-fortieth of the threshold, with a
horizon of 200 generations.

\begin{table}[h!]\centering
\begin{threeparttable}
\caption{Optimal clean-research share, relative to the threshold $R^\ast(0)$}
\label{tab:welfare_R}
\footnotesize
\begin{tabular}{llccccc}
\toprule
 & & & \multicolumn{2}{c}{Millian, actual tastes} & \multicolumn{2}{c}{Benthamite, actual tastes}\\
\cmidrule(lr){4-5}\cmidrule(lr){6-7}
Start & Scale effect & $R^\ast(0)$, \% & $\Psi=0$ & $\Psi=1$ & $\Psi=0$ & $\Psi=1$\\
\midrule
\multicolumn{7}{l}{\emph{A. Annual discount rate one per cent}}\\
Model baseline & $\nu=0$   & 1.59 & 0 & 1.63 & 3.05 & 3.83\\
               & $\nu=0.5$ & 1.99 & 0 & 1.40 & 2.55 & 3.17\\
Korea          & $\nu=0$   & 1.59 & 0 & 1.38 & 2.48 & 3.53\\
               & $\nu=0.5$ & 7.74 & 0 & 0.35 & 0.60 & 0.83\\
\midrule
\multicolumn{7}{l}{\emph{B. Annual discount rate a quarter of a per cent}}\\
Model baseline & $\nu=0$   & 1.59 & 0 & 1.15 & 2.15 & 2.32\\
               & $\nu=0.5$ & 1.99 & 0 & 1.08 & 1.85 & 2.00\\
Korea          & $\nu=0$   & 1.59 & 0 & 0.72 & 2.50 & 3.10\\
               & $\nu=0.5$ & 7.74 & 0 & 0.22 & 0.68 & 0.80\\
\bottomrule
\end{tabular}
\begin{tablenotes}\footnotesize
\item Optimal permanent clean-research share divided by the closed-form threshold $R^\ast(0)$ of
Corollary~\ref{cor:threshold}; values below one leave the economy below the threshold, so that population
declines indefinitely. No remediation; clean research financed by a labour-income tax; re-calibrated model
of Table~\ref{tab:recal}; flow utility of $4.87$ in 2023; horizon 200 generations.
\end{tablenotes}
\end{threeparttable}
\end{table}

Three results stand out. First, an average-utilitarian planner who places no direct value on the
environment chooses no clean research, whatever the start, the scale effect or the discount rate. The
cost falls on consumption and human capital in every generation, while the benefits, a population that
stops declining and a cleaner habitat, enter this criterion only through the small term $\alpha\ln b$.
With an amenity value, the same planner chooses a share above the threshold at the model baseline, by
between 8 and 63 per cent. Second, a total-utilitarian planner chooses between $1.9$ and $3.8$ times the
threshold at the model baseline: beyond securing a stationary population, extra clean research speeds the recovery of
habitat and so raises fertility and population along the way. Third, Korea at the upper end of the scale
effect is the exception. Its stationary population is a twenty-fifth of today's, so crossing the threshold
takes $7.7$ per cent of effective labour, and every criterion prefers to stay below it, at between a
fifth and four-fifths of the threshold, accepting an indefinite decline that begins centuries ahead.
Without a scale effect, Korea's threshold is the same as the baseline's, and the total-utilitarian
planner crosses it comfortably.

At the model baseline, a lower discount rate lowers the optimal share above the threshold. The gain from clean research beyond
the threshold comes early, as faster recovery, while its cost to income lasts forever, so giving the
future more weight favours less of it. A lower discount rate does not, in the range we consider, make the
Korean planner cross the threshold. The welfare case for clean research therefore rests on the value
placed on population, and, where fertility is lowest and the scale effect strong, the cost of avoiding
an environmental Empty Planet can exceed what any of these criteria would pay.

\clearpage
\section{Simulated country paths without a scale effect}
\label{app:paths}
\setcounter{table}{0}

Table~\ref{tab:proj} reports simulated paths for all nine countries at current clean progress and without
a scale effect, $\nu=0$, which is equivalent to continuing each country's 2005--2022 trend in emission
intensity. They illustrate the mechanism and should not be read as forecasts. For the seven countries that clear the threshold, the long-run levels coincide with the stationary
populations $L^\ast$ of Table~\ref{tab:thresholds}. Korea and Japan lie just below the threshold at
$\nu=0$ and do not settle; for them the table reports the population in 2383. The dates depend in addition on how quickly tastes and the
costs of children adjust to habitat, which the cross-sectional evidence does not identify, on the path of
income growth, which the simulation holds at its saturated rate, and on the omission of age structure and
migration. Under these assumptions fertility reaches $2.0$ by the mid-22nd century in France, the United
States and the United Kingdom, and only in the 23rd century in Italy and East Asia.

\begin{table}[h!]\centering
\begin{threeparttable}
\caption{Simulated population paths under the estimated habitat feedback}
\label{tab:proj}
\footnotesize
\setlength{\tabcolsep}{4pt}
\begin{tabular}{lcccccccc}
\toprule
 & TFR & $g_{\text{int}}$ & $\hat\zeta$ & $\epsilon^\ast$ & TFR $\ge2.0$ & \multicolumn{3}{c}{Population, 2023 $=1$}\\
\cmidrule(lr){7-9}
 & 2023 & \% a year & & & by & 2083 & 2143 & long run\\
\midrule
Korea          & 0.72 & $-1.67$ & 5.3  & 37.8 & 2293 & 0.18 & 0.09 & 0.04$^\dagger$\\
China          & 1.00 & $-2.71$ & 8.7  & 16.5 & 2263 & 0.30 & 0.18 & 0.10\\
Japan          & 1.21 & $-1.73$ & 5.5  & 9.2  & 2263 & 0.40 & 0.25 & 0.16$^\dagger$\\
Italy          & 1.20 & $-2.64$ & 8.4  & 9.4  & 2233 & 0.40 & 0.26 & 0.17\\
Sweden         & 1.43 & $-4.09$ & 13.1 & 5.1  & 2173 & 0.54 & 0.42 & 0.32\\
Germany        & 1.44 & $-3.18$ & 10.2 & 5.0  & 2173 & 0.54 & 0.41 & 0.32\\
United Kingdom & 1.56 & $-4.70$ & 15.0 & 3.7  & 2143 & 0.63 & 0.53 & 0.44\\
United States  & 1.62 & $-2.79$ & 8.9  & 3.1  & 2143 & 0.65 & 0.54 & 0.51\\
France         & 1.64 & $-3.08$ & 9.8  & 3.0  & 2143 & 0.67 & 0.57 & 0.54\\
\midrule
\multicolumn{9}{l}{No further decoupling (all countries): population declines indefinitely, by $0.41$ per cent a year in the limit}\\
\bottomrule
\end{tabular}
\begin{tablenotes}\footnotesize
\item Trend scenario: emission intensity $g_{\text{int}}$ continues at its 2005--2022 rate (CO$_2$ per unit
of GDP, Global Carbon Budget via Our World in Data); TFR from UN WPP 2024. $\epsilon^\ast$ is the habitat
per capita, relative to 2023, at which fertility reaches replacement. ``Long run'' is the stationary
population $Q_\infty/\epsilon^\ast$ of Proposition~\ref{prop:decoupling}(i). $\hat\zeta=-\zeta^\ast g_{\text{int}}/g_y$ expresses current clean progress as a
rate of decoupling, against the threshold $\zeta^\ast=5.88$. $^\dagger$Korea and Japan lie
just below the threshold; the entry is the 2383 value, after which population declines by
$0.04$ and $0.03$ per cent a year. Rows are ordered by long-run population.
\end{tablenotes}
\end{threeparttable}
\end{table}

\clearpage
\section{Cross-country evidence on the scale effect}
\label{app:nu}
\setcounter{table}{0}
The scale parameter $\nu$ governs how clean progress at a given research share varies with the number of
people. This appendix asks whether countries with more people decouple faster. By \eqref{eq:B_law},
clean progress per generation is $g_{B,i}=\delta_BR_iL_i^{\nu}$. If the clean-research share $R_i$ is
unrelated to country size once income is held fixed, the cross-country elasticity of $g_{B,i}$ with
respect to population identifies $\nu$.

We measure clean progress as in Section~\ref{sub:recal}, $g_{B,i}=(1+g_{\text{int},i})^{-30}-1$, where
$g_{\text{int},i}$ is the annual change in carbon-dioxide emissions per unit of GDP, for all countries with
more than one million people (Global Carbon Budget via Our World in Data). We report two specifications.
The first regresses $\ln g_{B,i}$ on log population, which requires $g_{B,i}>0$ and so drops the quarter of
countries whose emission intensity rose. The second keeps every country: it regresses the annual decline
in emission intensity, in levels, on log population and divides the coefficient by the mean decline. Both
control for log GDP per person, which absorbs structural change and the shift of emission-intensive
production to poorer countries, and for population growth over the period.

\begin{table}[h!]\centering
\begin{threeparttable}
\caption{Elasticity of clean progress with respect to population across countries}
\label{tab:nu}
\footnotesize
\begin{tabular}{llcccccc}
\toprule
 & & \multicolumn{3}{c}{(1) $\ln g_{B}$, countries with $g_B>0$} & \multicolumn{3}{c}{(2) decline in levels, all countries}\\
\cmidrule(lr){3-5}\cmidrule(lr){6-8}
Period & Sample & $\hat\nu$ & s.e. & $N$ & $\hat\nu$ & s.e. & $N$\\
\midrule
2005--2022 & All          & $-0.08$ & 0.06 & 110 & $0.09$  & 0.10 & 149\\
           & Richest 40\% & $-0.14$ & 0.09 & 58  & $-0.09$ & 0.05 & 60\\
1995--2022 & All          & $-0.10$ & 0.07 & 111 & $0.02$  & 0.08 & 147\\
           & Richest 40\% & $-0.26$ & 0.08 & 55  & $-0.12$ & 0.07 & 59\\
\bottomrule
\end{tabular}
\begin{tablenotes}\footnotesize
\item OLS across countries with more than one million people at the start of the period, controlling for
log GDP per person at the start and population growth over the period; heteroskedasticity-robust standard
errors. Column (2) divides the coefficient on log population by the mean annual decline in emission
intensity (for 2005--2022 and 1995--2022, 1.1 and 1.4 per cent a year for all countries, 2.8 and 2.6 for the
richest 40 per cent). Samples by
GDP per person at the start of the period.
\end{tablenotes}
\end{threeparttable}
\end{table}

Table~\ref{tab:nu} reports the results. No specification finds that larger countries decouple faster.
The point estimates lie between $-0.26$ and $0.09$, and the largest value inside any 95 per cent
confidence interval is $0.29$. A national scale effect of $0.5$ or more is therefore hard to reconcile with
the data.

Three qualifications limit what the test can show. First, research intensity may vary with size: small
countries may rely more on imported clean technology and do less clean research of their own, which would
bias the estimate upward, while large countries may be slower to reform, which would bias it downward.
Second, the decline in emission intensity reflects more than clean research, which is why we hold income
fixed; the same caveat applies to the calibration in Section~\ref{sub:recal}. Third, and most important,
clean knowledge crosses borders: clean innovation by firms draws on the clean knowledge produced in
the countries where they operate \citep{aghion2016carbon}. If clean knowledge is global, each country's
progress depends on world, not national, population, and a cross-country test cannot detect the scale
effect. In that case the relevant evidence is on research productivity over time, which falls as research
effort rises \citep{jones1995,bloom2020} and so argues against values near one.

The evidence therefore narrows the range worth considering without identifying $\nu$. We treat
$0\le\nu\le0.5$ as a sensitivity range: its lower end matches the cross-country estimates, and its upper end
exceeds every national estimate and leaves room for a global scale effect. The country thresholds in
Section~\ref{sub:proj} are reported at both ends, and at $\nu=1$.


\singlespacing
 {\bibliographystyle{chicago}
 	\bibliography{resources}
\end{document}
